% The next line makes it so it compiles the main file if I hit the compile button
% when editing this file in TeXworks.
% !TEX root = ../AIDMA.tex


\chapter{Graph Theory}\label{ch:graphs}

In this chapter we will provide a {\em very brief} and {\em very selective} introduction
to graphs.  Graph theory is a very wide field and there are many thick textbooks on the subject.
The main point of this chapter is to provide you
with the basic notion of what a graph is, some of the terminology used, a few applications,
and a few interesting and/or important results.

\section{Types of Graphs}\label{sec:graphTypes}
\begin{df}\index{graph}\index{simple graph}
\index{graph!simple}
\index{vertex}\index{edge}
A {\bf (simple) graph} $G$ = $(V,E)$ consists of 
\begin{itemize}
\setlength{\itemsep}{1pt}
\item
$V$, a nonempty set of {\bf vertices} and 
\item 
$E$, a set of {\em unordered pairs} of 
distinct vertices called {\bf edges}.
\end{itemize}
The {\bf order} of a graph is $|V|$, the number of vertices.
\end{df}

\begin{exa}
Here is an example of a graph with the set of vertices and edges listed on the right.
Vertices are usually represented by means of dots on the plane,
and the edges by means of lines connecting these dots.

\vspace*{-.1in}
\begin{center}
\sputfig{graphs/fig/lgraph.eps}{.55}
\end{center}
\end{exa}

\begin{exa}
Sometimes we just care about the visual representation of a graph.  Here are three examples.

\vspace*{-.25in}
\begin{center}
\sputfig{graphs/fig/simplegraph.eps}{.5}
\end{center}
\end{exa}


There are several variations of graphs.  
We will provide definitions and examples of the most common ones.

\begin{df}\index{directed graph}\index{digraph}\index{graph!directed}
A {\bf directed graph} (or {\bf digraph}) $G$ = $(V,E)$ consists of 
\begin{itemize}
\setlength{\itemsep}{1pt}
\item 
$V$, a nonempty set of {\bf vertices} and 
\item 
$E$, a set of {\em ordered pairs} of distinct vertices called {\bf directed edges}
(or just {\bf edges}).
\end{itemize}
The {\bf order} of a digraph is $|V|$, the number of vertices.
\end{df}

\begin{exa}
Here are three examples of directed graphs.
\begin{center}
%\sputfig{graphs/fig/digraph.eps}{.5} % Not necessary to have another example.
\sputfig{graphs/fig/directedgraph.eps}{.5}
\end{center}
\end{exa}

As you would probably suspect, the only difference between simple graphs and directed graphs
is that the edges in directed graphs have a direction.
We should note that simple graphs are sometimes called {\bf undirected graphs} to make it
clear that the graphs are not directed.

\begin{exa}
In a simple graph, $\{u,v\}$ and $\{v,u\}$ are just two different ways of talking about
the same edge--the edge between $u$ and $v$.  In a directed graph, $(u,v)$ is the edge from $u$ to $v$
and $(v,u)$ is the edge from $v$ to $u$. These are not the same, and they may or may not both be present.
\end{exa}


\begin{df}\index{multigraph}
A {\bf multigraph (directed multigraph)} $G=(V,E)$ consists of 
\begin{itemize}
\item
$V$, a set of vertices,
\item
$E$, a set of edges,  and
\item
a function $f$ from $E$ to 
{\small $\{\{u,v\}:u\neq v\in V\}$}

\noindent
(function $f$ from $E$ to 
{\small $\{(u,v):u\neq v\in V\}$.)}
\end{itemize}
Two edges $e_1$ and $e_2$ with 
$f(e_1)=f(e_2)$ are called {\bf multiple edges}.
\end{df}

Although the definition looks a bit complicated, 
a {\bf multigraph} $G=(V,E)$ is just a graph in which multiple
edges are allowed between a pair of vertices.

\begin{exa}
Here are a few examples of multigraphs.
\begin{center}
\sputfig{graphs/fig/multigraph.eps}{.55}
\end{center}
Here are some examples of directed multigraphs.
\begin{center}
\sputfig{graphs/fig/dmultigraph.eps}{.55}
\end{center}
\end{exa}



\begin{df}\index{pseudograph}\index{loop}
A {\bf pseudograph} $G=(V,E)$ is a graph in which we 
allow {\bf loops}--that is, edges from a vertex to itself.  
As you might imagine, a {\bf pseudo-multigraph} allows both loops
and multiple edges.
\end{df}

\begin{exa}
Here are some pseudographs.
\begin{center}
\sputfig{graphs/fig/psuedograph.eps}{.5}
\end{center}
Here are a few directed pseudographs.
\begin{center}
\sputfig{graphs/fig/dpsuedograph.eps}{.5}
\end{center}
\end{exa}


\begin{df}\index{weighted graph}\index{graph!weighted}
\index{network}
A {\bf weighted graph} is a graph (or digraph) with the additional
property that each edge $e$
has associated with it a real number $w(e)$ called its {\em weight}.

A {\bf weighted digraph} is often called a {\bf network}.
\end{df}

\begin{exa}
Here are two examples of weighted graphs and one weighted directed graph.
\begin{center}
\sputfig{graphs/fig/simplewgraph.eps}{.6}\hspace*{.5in}
\sputfig{graphs/fig/diwgraph.eps}{.4}
\end{center}
\end{exa}

As we have seen, there are several ways of categorizing graphs:
\begin{itemize}
\item
Directed or undirected edges.
\item
Weighted or unweighted edges.
\item
Allow multiple edges or not.
\item
Allow loops or not.
\end{itemize}

Unless specified, you can usually assume a graph does not allow
multiple edges or loops since these aren't that common.
Generally speaking, you can assume that if a graph is not specified as weighted or directed, 
it isn't.  The most common graphs we'll use are 
graphs, digraphs, weighted graphs, and networks.

\begin{rem}
When writing graph algorithms, it is important to know 
what characteristics the graphs have. For instance, if a graph might
have loops, the algorithm should be able to handle it.
Some algorithms do not work if a graph has loops and/or multiple edges,
and some only apply to directed (or undirected) graphs.  
\end{rem}

\newpage

\section{Graph Terminology}\label{sec:graphTerm}

\begin{df}
Given a graph $G=(V,E)$, we denote the number of vertices in $G$ by $|V|$ 
and the number of edges by $|E|$ (a notation that makes perfect sense since
$V$ and $E$ are sets).
\end{df}

\begin{df}\index{degree}\index{adjacent}\index{endpoint}\index{incident with}\index{connect}
Let $u$ and $v$ be vertices 
and $e=\{u,v\}$ be an edge 
in undirected graph $G$.
\begin{itemize}
\item
The vertices $u$ and $v$ are said to be {\bf adjacent}
\item
The vertices $u$ and $v$ are called the {\bf endpoints} of the edge $e$.
\item
The edge $e$ is said to be {\bf incident with} $u$ and $v$. 
\item
The edge $e$ is said to {\bf connect} $u$ and $v$.  
\item
The {\bf degree} of a vertex, denoted $deg(v)$, 
is the number of edges incident with it.  
\end{itemize}
\end{df}

\begin{exa}
Consider the following graphs.
 \begin{center}
\sputfig{graphs/fig/adjacent.eps}{.6}
\end{center}

\begin{minipage}[t]{.5\textwidth}
In graph $G_1$, we can say:
\begin{itemize}
\item
$w$ is adjacent to $x$.
\item
$w$ and $x$ are the endpoints of the edge $(w,x)$.
\item
$(w,x)$ is incident with both $w$ and $x$.
\item
$(w,x)$ connects vertices $w$ and $x$.
\end{itemize}
\end{minipage}
\hfill
\begin{minipage}[t]{.45\textwidth}
The following table gives the degree of each of the vertices in the graphs above.
\begin{center}
\begin{tabular}{|c|c|c|}
\hline
$G_1$&$G_2$&$G_3$\\
\hline
\hline
deg(u)=3&deg(u)=2&deg(u)=2\\
deg(v)=5&deg(v)=3&deg(v)=4\\
deg(w)=3&deg(w)=2&deg(w)=3\\
deg(x)=2&deg(x)=4&deg(x)=2\\
deg(y)=2&deg(y)=3&deg(y)=3\\
deg(z)=3&        &deg(z)=2\\
\hline
\end{tabular}
\end{center}
\end{minipage}
\end{exa}

\begin{df}
A {\bf subgraph} of a graph $G=(V,E)$ is a graph 
$G^{\prime}=(V^{\prime},E^{\prime})$ such that 
$V^{\prime}\subset V$ and 
$E^{\prime}\subset E$.
\end{df}

\begin{exa}
Consider the following three graphs:
\begin{center}
\sputfig{graphs/fig/subgraph.eps}{.35}
\end{center}
Notice that $H_2$ is a
subgraph of $H_1$ and that $H_3$ is a subgraph of both $H_1$ and $H_2$.
\end{exa}



\begin{df}\index{walk}
A $u-v$ {\bf walk} \index{walk} is an
alternating sequence of vertices and edges in $G$ with starting
vertex $u$ and ending vertex $v$ such that every edge joins the
vertices immediately preceding it and immediately following it.
\end{df}
You can think of a walk as follows: Put your pencil down on a vertex and
trace around edges however you like until you reach some destination vertex.
You are allowed to repeat edges and vertices as often as you like--just like you 
may repeat sidewalks and paths when you go for a walk (thus the name).

\begin{df}\index{trail}
A $u-v$ {\bf trail} \index{trail} is a $u-v$ walk that does not repeat an edge.
\end{df}

Notice that a trail {\em may} repeat a vertex.

\begin{df}\index{path}
A $u-v$ {\bf path}\index{path} is a walk that does not repeat any vertex.
\end{df}

It should be relatively easy to see that paths cannot repeat an edge
(because to repeat an edge you have to repeat a vertex).

\begin{exa}
In the first graph, the trail $abecde$ is indicated with the dark lines.
It is not a path since it repeats the vertex $e$.
The second and third graphs show examples of paths.
\begin{center}
\sputfig{graphs/fig/graph3.eps}{.55}\hspace*{.5in}
\sputfig{graphs/fig/graph4.eps}{.8}
\end{center}

Although not drawn (because it is harder to represent clearly on a drawing),
$acdecabecdeba$ is an example of a walk\footnote{This walk is
specified by just the vertices and not both the vertices and edges as in the definition. 
If multiple edges are not allowed (i.e. we are not working with a multigraph), then
there is no need to list the edges since they are clear.} 
To confirm it, you just need to verify
that there is an edge between adjacent vertices on the list.
On the other hand, $abcde$ is {\em not} a path, trail, or walk
because $(b,c)$ is not an edge.
\end{exa}

% DZ - simple path hasn't yet been defined
\begin{df}\index{cycle}\index{simple cycle}\index{graph!cycle}
A {\bf cycle} (or {\bf simple cycle}) 
is a list of vertices $v_1, v_2,\ldots, v_k$, with no repeats such that 
$(v_i,v_{i+1})$ is an edge for $i=1,\ldots, k-1$, and $(v_k,v_1)$ is an edge.

Put another way, a cycle is a path to which we append an edge from the last to the first vertex.

The number of vertices in a cycle is called its {\bf length}.
\end{df}

\begin{exa}\label{exa:cycle}
Here is a graph with a cycle of length $3$.
\begin{center}
\sputfig{graphs/fig/graph5.eps}{.4}
\end{center}
\end{exa}

\begin{exer}
Find a cycle of length 4 and a cycle of length 5 in the graph from Example~\ref{exa:cycle}.
Is there a cycle of length 6? Explain why or why not.

\noindent
\answerlinetwo
\begin{answer}
$abed$ is a cycle of length 4 and $ecdab$ is a cycle of length 5.
Other answers are possible.
There is no cycle of length 6 since there are only 5 vertices in the graph and a cycle cannot
repeat a vertex.
\end{answer}
\end{exer}

\begin{df}\index{connected, graph}\index{connected component}
A graph is called {\bf connected} if there is a path between every
pair of distinct vertices.

A {\bf connected component} of a graph is a maximal connected subgraph. 
\end{df}

\begin{exa}
Below are two graphs, each drawn inside dashed boxes.
The graph on the left is connected.  The one on the right
is not connected.  It has two connected components.
% The figures don't match and I don't currently have an easy way to edit them.
% CAC, 6/19/14
%The one on the right is also not connected, and it has three connected components.
\begin{center}
\sputfig{graphs/fig/graph6.eps}{.6}
%\hspace*{.25in}
%\sputfig{graphs/fig/graph7.eps}{.6}
\end{center}
\end{exa}


\begin{exer}
Draw a graph that has two connected components, one that is a cycle of length 4 and one
that is a cycle of length 3.

\vspace*{1in}
\begin{answer}
You should have drawn something like this (but probably bigger and with dots on the corners): $\square$   
$\triangle$
\end{answer}
\end{exer}


\begin{df}\index{tree}\index{unrooted tree}\index{forest}\index{leaf}
\begin{itemize}
\item
A {\bf tree} (or {\bf unrooted tree}) is a connected acyclic graph.
That is, a graph with no cycles.
\item
A vertex of degree one in a tree is called a {\bf leaf}.
\item
A {\bf forest} is a collection of trees.
\end{itemize}
\end{df}

\begin{exa}
Here are four trees.  If they were all part of the same graph, we could consider
the graph a forest.
\begin{center}
\putfig{graphs/fig/graph8.eps}
\end{center}
\end{exa}


\begin{exer}
Draw a tree that has 5 vertices, one vertex with degree 4 and the others with degree 1.

\vspace*{.5in}
\begin{answer}
You should have drawn something like this (but probably bigger and with dots on the corners 
and center): $\times$
\end{answer}
\end{exer}


\begin{exer}
Draw a forest with 5 trees that has 6 vertices.

\vspace*{.5in}
\begin{answer}
You should have drawn a path of length 2 and 4 vertices not connected to anything. Something like this:
 $|$ . . . .
\end{answer}
\end{exer}


\begin{rem}
These trees are not to be confused with {\em rooted trees} (e.g. {\em binary trees}).
When computer scientists use the term {\em tree}, they usually mean rooted trees,
not the trees we are discussing here.  When you see/hear the term `tree,' it is
important to be clear about which one the writer/speaker has in mind.
\end{rem}


\begin{df}\index{spanning tree}
A {\bf spanning tree} of $G$ is a subgraph which is a tree and
contains all of the vertices of $G$.
\end{df}

\begin{exa}
Below is a graph (on the left) and one of several possible
spanning trees (on the right).
\begin{center}
\putfig{graphs/fig/graph12.eps}
\end{center}
\end{exa}

% Don't need independent set for now, and have bipartite later.
%\begin{df}
%Let $G=(V,E)$ be a graph. A subset $S\subseteq V$ is an {\em
%independent set} of vertices if $uv\not\in E$ for all $u, v$ in
%$S$ ($S$ may be empty). A {\em bipartite graph}\index{bipartite graph} with bipartition
%$X, Y$ is a graph such that $V = X\cup Y$, $X\cap Y =
%\varnothing$, and $X$ and $Y$ are independent sets. $X$ and $Y$
%are called the {\em parts} of the bipartition.
%\end{df}


Here is some terminology related to directed graphs.

\begin{df}\index{adjacent to}\index{adjacent from}\index{initial vertex}\index{terminal vertex}\index{in-degree}\index{out-degree}\index{a graphs degreein@$deg^{-}$ (in-degree)}\index{a graphs degreeout@$deg^{+}$ (out-degree)}
Let $u$, $v$ be vertices in a directed graph $G$, and
$e=(u,v)$ be an edge in $G$.  
\begin{itemize}
\item $u$ is said to be {\bf adjacent to} $v$.
\item $v$ is said to be {\bf adjacent from} $u$.
\item $u$ is called the {\bf initial vertex} of $(u,v)$.
\item $v$ is called the {\bf terminal} or {\bf end vertex} of $(u,v)$.
\item
The {\bf in-degree} of $u$, denoted by $deg^{-}(u)$, 
is the number of edges in $G$ which have $u$ as their terminal vertex.
\item
The {\bf out-degree} of $u$, denoted by $deg^{+}(u)$, 
is the number of edges in $G$ which have $u$ as their initial vertex.
\end{itemize}
\end{df}

\newpage

\begin{exa}\label{exa:directedDegree}
Consider the three graphs below.
\begin{center}
\sputfig{graphs/fig/dadjacent.eps}{.5}
\end{center}
Consider the edge $(w,x)$ in $G_4$.
\begin{itemize}
\item
$w$ is adjacent to $x$ and $x$ is adjacent from $w$.
\item
$w$ is the initial vertex and $x$ is the terminal vertex of the edge $(w,x)$.
\end{itemize}
This table gives the in-degree and out-degree for the vertices in 
graphs $G_4$, $G_5$, and $G_6$.

\noindent
\begin{tabular}{|c|c||c|c||c|c|}
\hline
\multicolumn{2}{|c||}{$G_4$}&\multicolumn{2}{|c||}{$G_5$}&\multicolumn{2}{|c|}{$G_6$}\\
\hline
\hline
deg$^{-}$(u)=2&deg$^{+}$(u)=4&deg$^{-}$(u)=1&deg$^{+}$(u)=0&deg$^{-}$(u)=1&deg$^{+}$(u)=1\\
deg$^{-}$(v)=2&deg$^{+}$(v)=2&deg$^{-}$(v)=1&deg$^{+}$(v)=2&deg$^{-}$(v)=2&deg$^{+}$(v)=2\\
deg$^{-}$(w)=1&deg$^{+}$(w)=1&deg$^{-}$(w)=1&deg$^{+}$(w)=1&deg$^{-}$(w)=2&deg$^{+}$(w)=2\\
deg$^{-}$(x)=2&deg$^{+}$(x)=3&deg$^{-}$(x)=1&deg$^{+}$(x)=1&deg$^{-}$(x)=1&deg$^{+}$(x)=1\\
deg$^{-}$(y)=3&deg$^{+}$(y)=0&deg$^{-}$(y)=2&deg$^{+}$(y)=2&deg$^{-}$(y)=2&deg$^{+}$(y)=2\\
deg$^{-}$(z)=1&deg$^{+}$(z)=1&&&&\\
\hline
\end{tabular}
\end{exa}

\newpage

\section{Some Special Graphs}\label{sec:specialGraphs}

\begin{df}\index{complete graph}\index{graph!complete}
\index{a graphs complete@$K_n$ (complete graph)}
The {\bf complete graph} with $n$ vertices $K_n$ is the graph where 
every pair of vertices is adjacent. Thus, $K_n$ has $\binom{n}{2}$ edges.
\end{df}
\begin{exa}
Here are the complete graphs with $n=2,3,4,5$.
\begin{center}
\sputfig{graphs/fig/completegraphs.eps}{.45}
\end{center}
\end{exa}

\begin{df}\index{cycle}\index{graph!cycle}\index{a graphs cycle@$C_n$ (cycle)}
$C_n$ denotes a {\bf cycle}\index{cycle} of length $n$. It is a graph with $n$
edges, and $n$ vertices $v_1,\ldots , v_n$, where $v_i$ is adjacent
to $v_{i+1}$ for $n= 1, \ldots , n-1$, and $v_1$ is adjacent to $v_n$.
\end{df}

\begin{exa}\label{exa:cycles}
Here are the cycles of length $3$, $4$, and $5$.
\begin{center}
\sputfig{graphs/fig/cyclegraphs.eps}{.45}
\end{center}
\end{exa}

\begin{df}\index{path}\index{graph!path}\index{a graphs path@$P_n$ (path)}
$P_n$ denotes a {\em path}\index{path} of length $n$. It is a graph with $n$
edges, and $n+1$ vertices $v_0, v_1, \ldots , v_n$, where $v_i$ is
adjacent to $v_{i+1}$ for $n= 0, 1, \ldots , n-1$.
\end{df}

We won't provide an example of the paths because they are pretty easy to visualize.
For instance, $P_3$ is simply $C_4$ with one edge removed.

\begin{df}\index{graph!hypercube}\index{a graphs hypercube@$Q_n$ (hypercube)}
$Q_n$ denotes the {\bf $n$-dimensional cube} (or {\bf hypercube}).\index{hypercube} 
There are two equivalent ways to define $Q_n$. 
\begin{itemize}
\item
$Q_0$ is the graph consisting of a single vertex,
and $Q_n$ is constructed by taking two copies of $Q_{n-1}$
and connecting corresponding vertices between the two copies.
Hopefully you can see that this leads to 
$Q_1=P_2$ and $Q_2=C_4$. 
\item
$Q_n$ is the graph with 
$2^n$ vertices numbered $0$ through $2^n-1$
where two vertices are connected if the binary representation
of their numbers differs in exactly one place.
\end{itemize}
It is not difficult to determine that $Q_n$ has $n2^{n-1}$ edges.
\end{df}

\newpage
\begin{exa}
Here are $Q_2$ and $Q_3$, with vertices labeled as mentioned in the definition.
\begin{center}
\sputfig{graphs/fig/hypercube2}{1}
\hspace*{1in}
\sputfig{graphs/fig/hypercube3}{1}
\end{center}

Notice that in $Q_2$, the vertex labeled $11$ is adjacent to the vertices
labeled $10$ and $01$ since each of these differ in one bit.
Similarly, the vertex labeled $101$ in $Q_3$ is adjacent to the vertices
labeled $001$, $111$, and $100$ for the same reason. Next is $Q_4$, 
also labeled according to the definition.

\begin{center}
\sputfig{graphs/fig/hypercube4}{.75}
\end{center}

% The old versions--replaced by the ones above
% CAC, 6/1/2020
%
%\hfill
%\begin{minipage}{3.5cm}
%\begin{pspicture}(0,0)(1,.25)
%\rput(2,0){\begin{graph}(1,1) \roundnode{A}(-1,1)
%\roundnode{B}(1,1)\roundnode{C}(1,-1)\roundnode{D}(-1,-1)\edge{A}{B}\edge{B}{C}\edge{C}{D}\edge{D}{A}\autonodetext{A}[w]{01}
%\autonodetext{B}[e]{11}\autonodetext{C}[e]{10}\autonodetext{D}[w]{00}
%\end{graph}
%} 
%\rput(1.5,-1.95){$\mathbf{Q_2}$} 
%\end{pspicture}
%%\vspace*{.65in}\caption{$Q_2$.}\label{fig:q2}
%\end{minipage}
%\hfill
%\begin{minipage}{3.5cm}
%\begin{pspicture}(0,0)(1,.25)
%\rput(2,0){\begin{graph}(1,1) \roundnode{A}(-.5,.5)
%\roundnode{B}(.5,.5)\roundnode{C}(.5,-.5)\roundnode{D}(-.5,-.5)
% \roundnode{E}(-1.25,1.25)
%\roundnode{F}(1.25,1.25)\roundnode{G}(1.25,-1.25)\roundnode{H}(-1.25,-1.25)\edge{A}{B}\edge{B}{C}\edge{C}{D}
%\edge{D}{A}\autonodetext{A}[w]{011}
%\autonodetext{B}[e]{111}\autonodetext{C}[e]{101}\autonodetext{D}[w]{001}
%\autonodetext{E}[w]{010}\autonodetext{F}[e]{110}\autonodetext{G}[e]{100}\autonodetext{H}[w]{000}
%\edge{E}{F} \edge{F}{G}\edge{G}{H}\edge{H}{E}
%\edge{A}{E}\edge{B}{F}\edge{C}{G}\edge{D}{H}
%\end{graph}}
%\rput(1.5,-2.2){$\mathbf{Q_3}$} 
%\end{pspicture}
%%\vspace*{.65in}\caption{$Q_3$.}\label{fig:q3}
%\end{minipage}
%\hfill
%\vspace*{1in}

\end{exa}

% Replaced with the ones in the previous example that are better.
% CAC, 6/1/2020
%
%\begin{exa}
%Here are the (unlabeled) hypercubes with dimensions $1$, $2$, $3$, and $4$.
%\begin{center}
%\sputfig{graphs/fig/hypercube.eps}{.45}
%\end{center}
%\end{exa}

It should not be too difficult to see that $Q_1$ is the same as $P_1$ which is the same as $K_2$.

\begin{df}\index{bipartite graph}\index{graph!bipartite}
A simple graph $G$ is called {\bf bipartite} if
the vertex set $V$ can be partitioned into two disjoint nonempty sets $V_1$ and $V_2$
such that every edge connects a vertex in $V_1$ to a vertex in $V_2$.

Put another way, no vertices in $V_1$ are connected to each other, 
and no vertices in $V_2$ are connected to each other.
\end{df} 

Although there may be different ways of assigning the
vertices to $V_1$ and $V_2$, it does not matter.  
If there is at least one way to do so such that 
all edges go between $V_1$ and $V_2$, then $G$ is bipartite.

\begin{exa}\label{exa:bipart}
Here are a few bipartite graphs.  
\begin{center}
\sputfig{graphs/fig/bipartite.eps}{.5}
\end{center}Notice that although these are drawn to make it clear
what the partition is (i.e. $V_1$ is the top row of vertices and $V_2$ is the bottom row),
a graph does not have to be drawn as such in order to be bipartite.  
They are often drawn this way out of convenience.
For instance, the hypercubes are all bipartite even though they are not drawn this way.
\end{exa}

\begin{df}\index{complete bipartite graph}\index{graph!complete bipartite}
\index{a graphs bipartite@$K_{m,n}$ (complete bipartite graph)}
$K_{m,n}$ denotes the {\bf complete bipartite graph}\index{complete bipartite graph} with $m+n$
vertices. That is, it is the graph with $m+n$ vertices that is partitioned into two sets, 
one of size $n$ and the other of size $m$, such that every possible edge between the two sets is in the graph.
\end{df}

\begin{exa}
The first four graphs from Example~\ref{exa:bipart} are complete bipartite graphs.
The first is $K_{1,1}$, the second is $K_{1,2}$ (or $K_{2,1}$), the third is $K_{2,2}$, and the
fourth is $K_{3,2}$ (or $K_{2,3}$).
\end{exa}

\begin{exer}
Which of the following graphs are bipartite? 
Briefly justify your answers.
\begin{lenum}
\item $C_4$. \blanklinefill
\item $C_5$. \blanklinefill
\item $K_4$. \blanklinefill
\item $Q_3$. \blanklinefill
\item $P_n$ for any $n>0$. \blanklinefill
\end{lenum}
\begin{answer}
(a) Yes. Put the vertices on opposite corners in the same
 set.
(b) No. Every other vertex needs to be in the opposite
set, and since there is an odd cycle, eventually
we get to a situation where a vertex can't be in either set.
(c) No. None of the vertices can be in a set with 
any others since they are all connected to each other.
(d) Yes. If I put $000, 110, 011, 101$ in one set and
the others in the other set, you can easily
see that all of the edges go between sets.
(e) Yes. Since it is just a path, put every other vertex in 
the opposite set. Unlike the situation with the odd cycle,
we don't ever loop back so no problems can arise.
\end{answer}
\end{exer}


\begin{ques}
Prove or disprove: 
Every tree with at least 2 vertices is bipartite.
(Hint: You can prove this by induction on the number
of nodes.)
\vspace*{1.5in}
\begin{answer}
It is true. Notice that a tree with 2 vertices is clearly 
bipartite since it is just an edge.
Assume all trees with $n$ vertices are bipartite,
where $n\geq 2$.
Let $T$ be a tree with $n+1$ vertices and $v$ be a leaf
connected to some vertex $u$.
Let $T^{\prime}$ be the tree $T$ with the leaf $v$ removed. Then by in the inductive hypothesis, $T^{\prime}$ is bipartite since it has $n$ vertices.
Then $T$ is bipartite since we can put $v$ in the
opposite set of $u$ with no risk of edges between
the nodes within that set since
 $v$ is only connected to $u$.
Thus, by PMI, all trees with $n>2$ vertices are bipartite.
\end{answer}
\end{ques}
\newpage

\section{Handshaking Lemma}\label{sec:handshaking}
The following theorem is valid not only for simple graphs, 
but also for multigraphs and pseudographs.
\begin{thm}[Handshake Lemma]\label{thm:handshake_lemma}\index{handshake lemma}
Let $G=(V,E)$ be a  graph. Then $$ \sum _{v\in V} \deg(v) = 2|E|. $$
\begin{pf}
Let $X=\{(e,v) : e\in E, v\in V, \mbox{ and } e \mbox{ and } v \mbox{ are incident}\}$.
We will compute $|X|$ in two ways.
Each edge $e\in E$ is incident with exactly 2 vertices.  Thus, 
$$|X|=2 |E|.$$
Also, each vertex $v\in V$ is incident with deg($v$) edges. Thus, we
have that 
$${\displaystyle |X|=\sum_{v\in V} deg(v)}.$$
Setting these equal, we have the result.
\end{pf}
\end{thm}
The proof in the previous theorem is an example of a combinatorial proof. 
It is a neat technique where you prove a formula
by counting the number of objects in a set in two different ways. 

\begin{exa}
Consider the following graphs.
\begin{center}
\sputfig{graphs/fig/adjacent.eps}{.6}
\end{center}

A quick tabulation of the degrees of the vertices and the number of edges
reveals the following:
\begin{center}
\begin{tabular}{|l|c|c|c|}
\hline
Graph&$G_1$&$G_2$&$G_3$\\
\hline
\hline
$|E|$&9&7&8\\
\hline
${\displaystyle \sum_{v\in V} deg(v)}$&18&14&16\\
\hline
\end{tabular}
\end{center}
These results are certainly consistent with Theorem~\ref{thm:handshake_lemma}.
\end{exa}

Undirected graphs have an interesting property that is really easy to prove using 
Theorem~\ref{thm:handshake_lemma}.

\begin{cor}\label{cor:hand}
Every graph has an even number of vertices of odd degree.
\begin{pf}
The sum of an odd number of odd numbers is odd. Since the sum of
the degrees of the vertices in a simple graph is always even, one
cannot have an odd number of odd degree vertices.
\end{pf}
\end{cor}

The situation is slightly different, but not too surprising, for directed graphs.

\begin{thm}
Let $G=(V,E)$ be a directed graph.  Then
$$\sum_{v\in V} deg^{-}(v)=\sum_{v\in V} deg^{+}(v)=|E|.$$
\end{thm}


We won't provide a proof of this theorem (it's almost obvious), 
but you should verify it for the graphs in Example~\ref{exa:directedDegree} by
adding up the degrees in each column and comparing the appropriate sums.

\newpage

\section{Graph Representation}\label{sec:graphRep}

We provide only a very brief discussion of graph representation.  Consult your favorite data structure book for more details.

Let $G=(V,E)$ be a graph with $n$ vertices and $m$ edges. 
That is, $|V|=n$, and $|E|=m$.
There are two common ways of representing $G$. 
(There is actually a third, but it isn't nearly as common as the two we will discuss.)

The first method stores, for each vertex, a list of all of the vertices it is adjacent to.

\begin{df}\index{adjacency list}
The {\bf adjacency list} representation of a graph maintains, for each vertex, 
a list of all of the vertices adjacent to that vertex.
This can be implemented in many ways, but often an array of linked lists is used.
\end{df}

%Replaced with better examples. 6/2/2020, CAC.
%\begin{exa}
%The following shows a graph on the left with the adjacency list representation 
%on the right.
%\begin{center}
%\sputfig{graphs/fig/graph15.eps}{.65}
%\end{center}
%\end{exa}

\begin{exa}
A drawing of $C_5$ is given below on the left. 
An adjacency list representation is given below on the right.

\medskip
\noindent
\begin{minipage}{.3\textwidth}
 \rput(2,0){
 \begin{graph}(1,0)
 \roundnode{A}(1,0)
\roundnode{B}(.309,.951) %cos72,sin72
\roundnode{C}(-.809,.587) %cos144,sin144
\roundnode{D}(-.809,-.587) %cos216,sin216
\roundnode{E}(.309,-.951) %cos288,sin288
\autonodetext{A}[e]{A} \autonodetext{B}[n]{B}
\autonodetext{C}[w]{C} \autonodetext{D}[s]{D}
\autonodetext{E}[e]{E} \edge{A}{B} \edge{B}{C} \edge{C}{D}
\edge{D}{E} \edge{E}{A}
\end{graph}}
\end{minipage}
\hfill
\begin{minipage}{.2\textwidth}
\sputfig{graphs/fig/adjList1}{.75}
\end{minipage}
\hfill
\begin{minipage}{.4\textwidth}
Notice that every edge is represented twice.
For instance, the edge $(A,B)$ means that $A$ and $B$
are connected to each other. Thus, $B$ is on $A$'s list, 
and $A$ is on $B$'s list.
\end{minipage}
\end{exa}

\begin{exa}
A drawing of a directed cycle of length 5 is given below on the left. 
An adjacency list representation is given next to it.

\medskip
\noindent
\begin{minipage}{.3\textwidth}
 \rput(2,0){
 \begin{graph}(1,0)
 \roundnode{A}(1,0)
\roundnode{B}(.309,.951) %cos72,sin72
\roundnode{C}(-.809,.587) %cos144,sin144
\roundnode{D}(-.809,-.587) %cos216,sin216
\roundnode{E}(.309,-.951) %cos288,sin288
\autonodetext{A}[e]{A} \autonodetext{B}[n]{B}
\autonodetext{C}[w]{C} \autonodetext{D}[s]{D}
\autonodetext{E}[e]{E} \diredge{A}{B} \diredge{B}{C} \diredge{C}{D}
\diredge{D}{E} \diredge{E}{A}
\end{graph}}
\end{minipage}
\hfill
\begin{minipage}{.2\textwidth}
\sputfig{graphs/fig/adjList1b}{.75}
\end{minipage}
\hfill
\begin{minipage}{.4\textwidth}
Notice that this is a lot like the previous example except that each list only has one element on it.
That is because $(A,B)$ is an edge (for instance), but $(B,A)$ is not an edge. So $B$ is on $A$'s list, but
$A$ is not on $B$'s list.
\end{minipage}
\end{exa}

\begin{exa}\label{exa:thirdAL}
Here is another example of a graph on the left with the adjacency list representation on the right.

\medskip
\noindent
\begin{minipage}{.3\textwidth}
 \rput(2,0){\begin{graph}(1,0)
\roundnode{A}(-1,1) 
 \roundnode{B}(1,1)
  \roundnode{D}(-1,-1)
  \roundnode{C}(1,-1)
  \roundnode{E}(0,0)
\autonodetext{A}[n]{$A$} 
\autonodetext{B}[n]{$B$} 
\autonodetext{C}[s]{$C$} 
\autonodetext{D}[s]{$D$}
\autonodetext{E}[w]{$E$} 
\edge{A}{B}
\edge{A}{E} 
\edge{B}{C}
\edge{B}{E} 
\edge{C}{D} 
\edge{C}{E} 
\edge{D}{A}
\edge{D}{E} 
 \end{graph}}
 \end{minipage}
\hfill
\begin{minipage}{.6\textwidth}
\sputfig{graphs/fig/adjList2}{.75}
\end{minipage}

%Don't know why this is needed but it's too squished without it.
\vspace*{.125in}

\noindent
Note that the order the vertices are listed does not matter.
\end{exa}


\begin{exer}
Give the adjacency list representation for $K_{3,3}$ as drawn below.

\noindent
\begin{minipage}{.3\textwidth}
 \rput(2,0){\begin{graph}(1,1)
\roundnode{A}(-1,1) \roundnode{B}(0,1)
 \roundnode{C}(1,1)
  \roundnode{D}(-1,-1)
  \roundnode{E}(0,-1)
  \roundnode{F}(1,-1)
\autonodetext{A}[n]{$A$} \autonodetext{B}[n]{$B$}
\autonodetext{C}[n]{$C$} \autonodetext{D}[s]{$D$}
\autonodetext{E}[s]{$E$} \autonodetext{F}[s]{$F$} \edge{A}{D}
\edge{A}{E}\edge{A}{F} \edge{B}{D} \edge{B}{E}\edge{B}{F}
\edge{C}{D} \edge{C}{E}\edge{C}{F}
 \end{graph}}
 \end{minipage}
 \hfill
\begin{minipage}{.6\textwidth}
\sputfig{graphs/fig/adjListEmpty}{1}
\end{minipage}
\begin{answer}

\sputfig{graphs/fig/adjListK33}{.75}
\end{answer}
\end{exer}

\begin{exer}
Give the adjacency list representation for the directed graph similar to $K_{3,3}$ drawn below.

\noindent
\begin{minipage}{.3\textwidth}
 \rput(2,0){\begin{graph}(1,1)
\roundnode{A}(-1,1) \roundnode{B}(0,1)
 \roundnode{C}(1,1)
  \roundnode{D}(-1,-1)
  \roundnode{E}(0,-1)
  \roundnode{F}(1,-1)
\autonodetext{A}[n]{$A$} \autonodetext{B}[n]{$B$}
\autonodetext{C}[n]{$C$} \autonodetext{D}[s]{$D$}
\autonodetext{E}[s]{$E$} \autonodetext{F}[s]{$F$} \diredge{A}{D}
\diredge{A}{E}\diredge{A}{F} \diredge{B}{D} \diredge{B}{E}\diredge{B}{F}
\diredge{C}{D} \diredge{C}{E}\diredge{C}{F}
 \end{graph}}
 \end{minipage}
 \hfill
\begin{minipage}{.6\textwidth}
\sputfig{graphs/fig/adjListEmpty}{1}
\end{minipage}
\begin{answer}

\sputfig{graphs/fig/adjListK33b}{.75}
\end{answer}
\end{exer}

The graph in Example~\ref{exa:thirdAL} has 5 vertices and 8 edges (so $n=5$ and $m=8$).
The adjacency list uses an array of size 5 and there are 5 linked lists that contain 
a total of $3+3+3+3+4=16=2*8$ nodes. Notice that this is twice the number of edges because
each edge is stored twice (because if $(u,v)$ is an edge, $u$ is stored on $v$'s list and $v$ is stored on $u$'s list).
For each node we need to store the {\em value} and the {\em next} node,
so the linked lists take up about $2(2*8)=4*8=4m$ memory.
Since the array takes about $5=n$ memory, the memory requirement for an adjacency list representation of the graph is approximately $n+4m=\Theta(n+m)$. 

Notice that the discussion in the previous paragraph generalizes to all graphs.
That is, the space requirement for the adjacency list representation of a graph is approximately $n+4m=\Theta(n+m)$.

Hopefully it is not too difficult to see that for directed graphs, the amount of memory required is
about $n+2m=\Theta(n+m)$ because each edge is only stored once.

For weighted graphs, an additional field can be stored in each node for the weight of each edge. 
So for undirected weighted graphs, the memory requirement goes up to about $n+6m$,
and for directed weighted graphs it is about $n+3m$. In both cases, it is still $\Theta(n+m)$.

The second method of storing a graph makes it so you can ask directly ``Is $(u,v)$ and edge?''
This is accomplished by storing a matrix whose rows and columns are indexed by the vertices.

\newpage % bad break.

\begin{df}\index{adjacency matrix}
The {\bf adjacency matrix} $M$ of a graph $G$ is the
$n$ by $n$ matrix $M$ defined as
$$M(i,j)=\left\{
\begin{array}{ll}
1&\mbox{ if } (i,j) \mbox{ is an edge}\\
0&\mbox{ if } (i,j) \mbox{ is not an edge}\\
\end{array}
\right.$$
\end{df}
We often assume that the vertices are numbered $0,1,\ldots, n-1$
since that is how we typically index matrices. 
In the next few examples we will continue with our examples with vertices labeled $A$, $B$, etc.
To make the interpretation of the matrices clear, we label the rows and columns.
You can also just think of a mapping of $A$ to 0, $B$ to 1, etc.

\begin{exa}
A drawing of $C_5$ is given below on the left, the adjacency list in the middle, 
and the adjacency matrix on the right.

\medskip
\noindent
\begin{minipage}{.30\textwidth}
 \rput(2,0){
 \begin{graph}(1,0)
 \roundnode{A}(1,0)
\roundnode{B}(.309,.951) %cos72,sin72
\roundnode{C}(-.809,.587) %cos144,sin144
\roundnode{D}(-.809,-.587) %cos216,sin216
\roundnode{E}(.309,-.951) %cos288,sin288
\autonodetext{A}[e]{A} \autonodetext{B}[n]{B}
\autonodetext{C}[w]{C} \autonodetext{D}[s]{D}
\autonodetext{E}[e]{E} \edge{A}{B} \edge{B}{C} \edge{C}{D}
\edge{D}{E} \edge{E}{A}
\end{graph}}
\end{minipage}
\hfill
\begin{minipage}{.3\textwidth}
\sputfig{graphs/fig/adjList1}{.75}
\end{minipage}
\hfill
\begin{minipage}{.3\textwidth}
$
\begin{blockarray}{*{6}{c}}
 &A & B & C & D & E \\
\begin{block}{c[*{5}{c}]}
A&0&1&0&0&1  \bigstrut[t]\\
B&1&0&1&0&0\\
C&0&1&0&1&0\\
D&0&0&1&0&1\\
E&1&0&0&1&0\\
\end{block}
\end{blockarray}
$
\end{minipage}
\end{exa}

\begin{exa}\label{exa:digraph1}
A drawing of a directed cycle of length 5 is given below on the left. 
An adjacency list representation is given in the middle and 
the adjacency matrix on the right.

\medskip
\noindent
\begin{minipage}{.3\textwidth}
 \rput(2,0){
 \begin{graph}(1,0)
 \roundnode{A}(1,0)
\roundnode{B}(.309,.951) %cos72,sin72
\roundnode{C}(-.809,.587) %cos144,sin144
\roundnode{D}(-.809,-.587) %cos216,sin216
\roundnode{E}(.309,-.951) %cos288,sin288
\autonodetext{A}[e]{A} \autonodetext{B}[n]{B}
\autonodetext{C}[w]{C} \autonodetext{D}[s]{D}
\autonodetext{E}[e]{E} \diredge{A}{B} \diredge{B}{C} \diredge{C}{D}
\diredge{D}{E} \diredge{E}{A}
\end{graph}}
\end{minipage}
\hfill
\begin{minipage}{.3\textwidth}
\sputfig{graphs/fig/adjList1b}{.75}
\end{minipage}
\hfill
\hfill
\begin{minipage}{.3\textwidth}
$
\begin{blockarray}{*{6}{c}}
 &A & B & C & D & E \\
\begin{block}{c[*{5}{c}]}
A&0&1&0&0&0  \bigstrut[t]\\
B&0&0&1&0&0\\
C&0&0&0&1&0\\
D&0&0&0&0&1\\
E&1&0&0&0&0\\
\end{block}
\end{blockarray}
$
\end{minipage}
\end{exa}


\begin{exa}%Oops. Duplicate. \label{exa:thirdAL}
Here is another example of a graph on the left,
the adjacency list representation on the center,
and the adjacency matrix on the right.

\medskip
\noindent
\begin{minipage}{.3\textwidth}
 \rput(2,0){\begin{graph}(1,0)
\roundnode{A}(-1,1) 
 \roundnode{B}(1,1)
  \roundnode{D}(-1,-1)
  \roundnode{C}(1,-1)
  \roundnode{E}(0,0)
\autonodetext{A}[n]{$A$} 
\autonodetext{B}[n]{$B$} 
\autonodetext{C}[s]{$C$} 
\autonodetext{D}[s]{$D$}
\autonodetext{E}[w]{$E$} 
\edge{A}{B}
\edge{A}{E} 
\edge{B}{C}
\edge{B}{E} 
\edge{C}{D} 
\edge{C}{E} 
\edge{D}{A}
\edge{D}{E} 
 \end{graph}}
 \end{minipage}
\hfill
\begin{minipage}{.3\textwidth}
\sputfig{graphs/fig/adjList2}{.75}
\end{minipage}
\hfill
\begin{minipage}{.3\textwidth}
$
\begin{blockarray}{*{6}{c}}
 &A & B & C & D & E \\
\begin{block}{c[*{5}{c}]}
A&0&1&0&1&1  \bigstrut[t]\\
B&1&0&1&0&1\\
C&0&1&0&1&1\\
D&1&0&1&0&1\\
E&1&1&1&1&0\\
\end{block}
\end{blockarray}
$
\end{minipage}
\end{exa}

\begin{rem}
You may or may not have noticed from these examples
that the adjacency matrix of an undirected graph is always
symmetric. If you think about it, the reason should be obvious.

The same cannot be said of directed graphs, as should be obvious from Example~\ref{exa:digraph1}.
\end{rem}

From these examples, it should be relatively clear that the amount of space
needed to store an adjacency matrix with $n$ vertices and $m$ edges is about $n^2=\Theta(n^2)$.
Notice that it does not depend on $m$, since a larger $m$ just means more 1s and 
fewer 0s in the matrix.

If $G$ is weighted, we can store the weights in the matrix instead of just $0$ or $1$.
For non-adjacent vertices, we store $\infty$, or {\em  MAX\_INT}
(or $-1$ if only positive weights are valid).  If done this way, the space requirement remains
$n^2=\Theta(n^2)$.
Alternatively, a second matrix can be used to store the weights, doubling the space requirement,
which is still $\Theta(n^2)$.

Notice the amont of space required to store both directed and undirected graphs is the same
with the adjacency matrix.

\begin{exer}
Give the adjacency matrix representation for $K_{3,3}$ as drawn below.

\bigskip
\bigskip
\noindent
 \rput(2,0){\begin{graph}(1,1)
\roundnode{A}(-1,1) \roundnode{B}(0,1)
 \roundnode{C}(1,1)
  \roundnode{D}(-1,-1)
  \roundnode{E}(0,-1)
  \roundnode{F}(1,-1)
\autonodetext{A}[n]{$A$} \autonodetext{B}[n]{$B$}
\autonodetext{C}[n]{$C$} \autonodetext{D}[s]{$D$}
\autonodetext{E}[s]{$E$} \autonodetext{F}[s]{$F$} \edge{A}{D}
\edge{A}{E}\edge{A}{F} \edge{B}{D} \edge{B}{E}\edge{B}{F}
\edge{C}{D} \edge{C}{E}\edge{C}{F}
 \end{graph}}

%Aparently the stupid drawing doesn't have a height? Whatever.
\vspace*{1.5in}

\begin{answer}

$
\begin{blockarray}{*{7}{c}}
 &A & B & C & D & E &E \\
\begin{block}{c[*{6}{c}]}
A&0&0&0&1&1&1  \bigstrut[t]\\
B&0&0&0&1&1&1\\
C&0&0&0&1&1&1\\
D&1&1&1&0&0&0\\
E&1&1&1&0&0&0\\
F&1&1&1&0&0&0\\
\end{block}
\end{blockarray}
$
\end{answer}
\end{exer}

\begin{exer}
Give the adjacency matrix representation for the directed graph similar to $K_{3,3}$ drawn below.

\bigskip
\bigskip
\noindent
 \rput(2,0){\begin{graph}(1,1)
\roundnode{A}(-1,1) \roundnode{B}(0,1)
 \roundnode{C}(1,1)
  \roundnode{D}(-1,-1)
  \roundnode{E}(0,-1)
  \roundnode{F}(1,-1)
\autonodetext{A}[n]{$A$} \autonodetext{B}[n]{$B$}
\autonodetext{C}[n]{$C$} \autonodetext{D}[s]{$D$}
\autonodetext{E}[s]{$E$} \autonodetext{F}[s]{$F$} \diredge{A}{D}
\diredge{A}{E}\diredge{A}{F} \diredge{B}{D} \diredge{B}{E}\diredge{B}{F}
\diredge{C}{D} \diredge{C}{E}\diredge{C}{F}
 \end{graph}}
 
%Aparently the stupid drawing doesn't have a height? Whatever.
\vspace*{1.5in}

\begin{answer}

$
\begin{blockarray}{*{7}{c}}
 &A & B & C & D & E &E \\
\begin{block}{c[*{6}{c}]}
A&0&0&0&1&1&1  \bigstrut[t]\\
B&0&0&0&1&1&1\\
C&0&0&0&1&1&1\\
D&0&0&0&0&0&0\\
E&0&0&0&0&0&0\\
F&0&0&0&0&0&0\\
\end{block}
\end{blockarray}
$
\end{answer}
\end{exer}

Obviously, how much space is required to store a graph is of importance, but so is how much
time is required to do basic operations on a graph. For instance, the most common things
one might want to do on a graph are determine whether or not two vertices are adjacent
and iterate over the edges that are incident with a vertex (put another way, iterate
over all of the neighbors of a vertex). For a weighted graph, one would probably ask
the weight of an edge somewhat often. There are certainly other important operations
one might want to perform on a graph.
Since you have all of the tools you need to 
answer such questions, we will ask  you to explore them at the end of the chapter.

So which representation is better? We will also let you think about that at the end of the chapter,
but hopefully it is somewhat clear that answering that question requires you to consider both
time and space requirements.

\newpage

\section{Problem Solving with Graphs}\label{sec:GraphProbSolve}

There are many problems on graphs that are of interest for various reasons.
The following very short list contains some of the more common ones.

\begin{itemize}
\item  {\sc Path}: Is there a path from A to B?
\item {\sc Cycles}:
  Does the graph contain a cycle?
\item {\sc Connectivity}:
  Is there a way to get between any two vertices in the graph?
\item {\sc Biconnectivity}:
  Will the graph become disconnected if
  one vertex is removed?
\item {\sc Planarity}:
  Is there a way to draw the graph without
  edges crossing?
\item {\sc Shortest Path}:
  What is the shortest path from A to B? (weighted and unweighted versions)
\item {\sc Longest Path}:
  What is the longest path from A to B? (weighted and unweighted versions)
\item {\sc Minimum Spanning Tree}:
  What is the ``most efficient'' way to connect the vertices (weighted graphs)?
  \item {\sc Traversability}:
  Is is possible to travel to every vertex without repeating a vertex?
  Is it possible to travel over every edge without repeating an edge?
\item {\sc Traveling Salesman}:
  What is the shortest route that visits every vertex and returns to the starting vertex? 
  (weighted graphs)
\end{itemize}

Knowing what graph problems have been studied and what is known about each is
very important.  Many problems can be modeled using graphs, and once a problem
has been mapped to a particular graph problem, it can be helpful to know the best 
way to solve it.

We finish the chapter by giving several examples of problems whose solutions become
simpler when using a graph-theoretic model as well as develop some new
graph terminology. 
It is important to mention that whole books are written just about graph theory, 
and even they have to pick a small subset of the topic. Thus, what is presented in the
remainder of this chapter should not be interpreted in any way to be the most important
topics in graph theory. It is just a very small selection of easy to understand topics
that are related to interesting problems. Dozens--maybe even hundreds--of other topics
could have been chosen.
It should be noted that the author has even resisted the urge to include 
one of his favorite graph topics, {\em graph pebbling}, even though it is a somewhat
interesting topic. Well, to him anyway.


\newpage

\subsection{Sample Problems}
\begin{exa}\label{exa:wolf_goat_cabbage}\index{wolf}\index{cabbage}\index{goat}\index{boat}\index{ferryman}
A wolf, a goat, and a cabbage are on one bank of a river. The
ferryman wants to take them across, but his boat is too small to
accommodate more than one of them at a time. 
He cannot leave the wolf and the goat together (the wolf will eat the goat), 
or the cabbage and the goat (the goat will eat the cabbage) unless he is with them. 
Can the ferryman still get all of them across the river?
\begin{sol}
Represent the position of a single item by $0$ for one bank of the 
river and  $1$ for the other bank. The position
of the three items can now be given as an ordered triplet, say
$(W,G,C)$. For example, $(0,0,0)$ means that the three items are
on one bank of the river, $(1,0,0)$ means that the wolf is on one
bank of the river while the goat and the cabbage are on the other
bank. The object of the puzzle is now seen to be to move from
$(0,0,0)$ to $(1,1,1)$ by traversing certain edges of $Q_3$ while avoiding
other edges. Note that $Q_3$ is the correct set of edges to consider 
since he can only move one of the three items at a time.

But there are some edges he cannot use.
For instance, $000\rightarrow 100$ is illegal since
it would mean he takes the wolf to the other side, leaving the goat and cabbage together.
Similarly, $000\rightarrow 001$ is illegal. Thus, from $000$, the only choice is to go to $010$.
Continuing this analysis, it can be determined that the set of legal edges is as in the following graph:

%can be seen in figure  \ref{fig:wolf_goat_cabbage}.
%\begin{figure}[h]
\begin{center}
\begin{graph}(1,1)
\roundnode{A}(1,0) \roundnode{B}(.5,.866) \roundnode{C}(-.5,.866)
\roundnode{D}(-1,0) \roundnode{E}(-.5,-.866)
\roundnode{F}(.5,-.866) \roundnode{G}(-2,0) \roundnode{H}(2,0)
\autonodetext{A}[w]{101} \autonodetext{B}[n]{001}
\autonodetext{C}[n]{011} \autonodetext{D}[e]{010}
\autonodetext{E}[s]{110} \autonodetext{F}[s]{100}
\autonodetext{G}[w]{000}\autonodetext{H}[e]{111} \edge{H}{A}
\edge{A}{B} \edge{B}{C} \edge{C}{D} \edge{D}{E} \edge{E}{F}
\edge{F}{A}\edge{G}{D}
\end{graph}\vspace*{.5in}
\end{center}
%\caption{Example\ref{exa:wolf_goat_cabbage}.}\label{fig:wolf_goat_cabbage}
%\end{figure}

Based on this, one answer is  $000\rightarrow 010\rightarrow 011
\rightarrow 001\rightarrow 101 \rightarrow 111.$
This means that
the ferryman (i) takes the goat across, (ii) returns and takes the
cabbage over, (iii) brings back the goat, (iv) takes the wolf over,
(v) returns and takes the goat over. 

Another answer is $000\rightarrow 010\rightarrow 110 \rightarrow
100\rightarrow 101 \rightarrow 111.$
This means that the ferryman
(i) takes the goat across, (ii) returns and takes the wolf over,
(iii) brings back the goat, (iv) takes the cabbage over, (v) returns
and takes the goat over.

Go to \url{https://xkcd.com/1134/} to see a funny, but incorrect, solution.
\end{sol}
\end{exa}


\begin{exa}
%[E\"{o}tv\"{o}s Mathematical Competition, 1947] 
\label{exa:ramsey6-2}Prove that amongst
six people in a room there are at least three who know one
another, or at least three who do not know one another.
\begin{sol}
%In graph-theoretic terms, we need to show that every colouring of the 
%edges of $K_6$ into two different colours, say red and blue, contains
%a monochromatic  triangle (that is, the edges of the triangle have
%all the same colour). 
Consider an arbitrary person of this group
(call him Peter). There are five other people, and of these,
either three of them know Peter or else, three of them do not know
Peter. 

Let us assume three know Peter.
If two of these three people know one another,
then we have a triangle of three people who know each other 
(Peter and these two--see the graph below on the left, 
where the acquaintances are marked by solid
lines). If no two of these three people know one another, then we
have three mutual strangers (see the graph on the right).

%\begin{figure}[h]
\begin{minipage}{.4\textwidth}
\centering
\begin{graph}(1,1) \roundnode{A}(-1,0) \roundnode{B}(1,1)
\roundnode{C}(2,.7) \roundnode{D}(1.5,.2) \roundnode{E}(2,-.7)
\roundnode{F}(1,-1)\autonodetext{A}[w]{Peter}  \edge{A}{B}
\edge{A}{C} \edge{A}{D} \edge{B}{D} \edge{A}{E}[\graphlinedash{3}]
\edge{A}{F}[\graphlinedash{3}]
\edge{B}{C}[\graphlinedash{3}]\edge{C}{D}[\graphlinedash{3}]
\end{graph}\vspace*{.5in}
%\caption{Example \ref{exa:ramsey6-2}.}\label{fig:ramsey6-2}
\end{minipage}
\hfill
\begin{minipage}{.4\textwidth}
\centering
\begin{graph}(1,1) \roundnode{A}(-1,0) \roundnode{B}(1,1)
\roundnode{C}(2,.7) \roundnode{D}(1.5,.2) \roundnode{E}(2,-.7)
\roundnode{F}(1,-1)\autonodetext{A}[w]{Peter}  \edge{A}{B}
\edge{A}{C} \edge{A}{D} \edge{B}{D}[\graphlinedash{3}]
\edge{A}{E}[\graphlinedash{3}] \edge{A}{F}[\graphlinedash{3}]
\edge{B}{C}[\graphlinedash{3}]\edge{C}{D}[\graphlinedash{3}]
\end{graph}\vspace*{.5in}
%\caption{Example \ref{exa:ramsey6-2}.}\label{fig:ramsey6-2_2}
\end{minipage}
%\end{figure}

The argument for the case when three do not know Peter is similar
and is left to the reader.

\end{sol}
\end{exa}



\begin{exa}\label{exa:landaus}Mr. and Mrs. Landau invite four other married couples
for dinner. Some people shook hands with some others, and the
following rules were noted: (i) a person did not shake hands with
himself, (ii) no one shook hands with his spouse, (iii) no one
shook hands more than once with the same person. After the
introductions, Mr. Landau asks the nine people how many hands they
shook. Each of the nine people asked gives a different number. How
many hands did Mrs. Landau shake?
\begin{sol}
The given numbers can either be $0,1,2,\ldots , 8$, or
$1,2,\ldots , 9$. Now, the sequence $1,2,\ldots , 9$ must be ruled
out, since if a person shook hands nine times, then he must have
shaken hands with his spouse, which is not allowed. The only
permissible sequence is thus $0,1,2,\ldots , 8$. Consider the
person who shook hands $8$ times, as in figure \ref{fig:landaus1}.
Discounting himself and his spouse, he must have shaken hands with
everybody else. This means that he is married to the person who
shook $0$ hands! We now consider the person that shook $7$ hands,
as in figure \ref{fig:landaus2}. He didn't shake hands with
himself, his spouse, or with the person that shook $0$ hands. But
the person that shook hands only once did so with the person
shaking $8$ hands. Thus the person that shook hands $7$ times is
married to the person that shook hands once. Continuing this
argument, we see the following pairs: $(8,0)$, $(7,1)$, $(6,2)$,
$(5,3)$. This leaves the person that shook hands $4$ times without
a partner, meaning that this person's partner did not give a
number, hence this person must be Mrs. Landau! Conclusion: Mrs.
Landau shook hands four times. A graph of the situation appears in
figure \ref{fig:landaus3}. \vspace{1cm}
\end{sol}
\end{exa}
\begin{figure}[ht]
\begin{minipage}{.3\textwidth}
\centering
\begin{graph}(1,1)
\roundnode{A}(1,0)
\roundnode{B}(.809, .588) %cos36,sin36
\roundnode{C}(.309, .951) %cos72,sin72
\roundnode{D}(-.309, .951) %cos108,sin108
\roundnode{E}(-.809, .588) %cos144,sin144
\roundnode{F}(-1, 0) %cos180,sin180
\roundnode{G}(-.809, -.588) %cos216,sin216
\roundnode{H}(-.309, -.951) %cos252,sin252
\roundnode{I}(.309, -.951) %cos288,sin288
\roundnode{J}(.809, -.588) %cos324,sin324
\autonodetext{A}[e]{Mr. Landau} \autonodetext{B}[e]{8}
\autonodetext{C}[n]{7} \autonodetext{D}[n]{6}
\autonodetext{E}[n]{5} \autonodetext{F}[w]{4}
\autonodetext{G}[s]{3} \autonodetext{H}[s]{2}
\autonodetext{I}[s]{1} \autonodetext{J}[s]{0} \edge{B}{A}
\edge{B}{C} \edge{B}{D} \edge{B}{E} \edge{B}{F} \edge{B}{G}
\edge{B}{H} \edge{B}{I}
\end{graph}\vspace*{.5in}\caption{Example
\ref{exa:landaus}.} \label{fig:landaus1}
\end{minipage}
\hfill
\begin{minipage}{.3\textwidth}
\centering
\begin{graph}(1,1)
\roundnode{A}(1,0)
\roundnode{B}(.809, .588) %cos36,sin36
\roundnode{C}(.309, .951) %cos72,sin72
\roundnode{D}(-.309, .951) %cos108,sin108
\roundnode{E}(-.809, .588) %cos144,sin144
\roundnode{F}(-1, 0) %cos180,sin180
\roundnode{G}(-.809, -.588) %cos216,sin216
\roundnode{H}(-.309, -.951) %cos252,sin252
\roundnode{I}(.309, -.951) %cos288,sin288
\roundnode{J}(.809, -.588) %cos324,sin324
\autonodetext{A}[e]{Mr. Landau} \autonodetext{B}[e]{8}
\autonodetext{C}[n]{7} \autonodetext{D}[n]{6}
\autonodetext{E}[n]{5} \autonodetext{F}[w]{4}
\autonodetext{G}[s]{3} \autonodetext{H}[s]{2}
\autonodetext{I}[s]{1} \autonodetext{J}[s]{0} \edge{B}{A}
\edge{B}{C} \edge{B}{D} \edge{B}{E} \edge{B}{F} \edge{B}{G}
\edge{B}{H} \edge{B}{I} \edge{C}{D} \edge{C}{E} \edge{C}{F}
\edge{C}{G}\edge{C}{H}\edge{C}{A}
\end{graph}\vspace*{.5in}\caption{Example
\ref{exa:landaus}.} \label{fig:landaus2}
\end{minipage}
\hfill
\begin{minipage}{.3\textwidth}
\centering
\begin{graph}(1,1)
\roundnode{A}(1,0)
\roundnode{B}(.809, .588) %cos36,sin36
\roundnode{C}(.309, .951) %cos72,sin72
\roundnode{D}(-.309, .951) %cos108,sin108
\roundnode{E}(-.809, .588) %cos144,sin144
\roundnode{F}(-1, 0) %cos180,sin180
\roundnode{G}(-.809, -.588) %cos216,sin216
\roundnode{H}(-.309, -.951) %cos252,sin252
\roundnode{I}(.309, -.951) %cos288,sin288
\roundnode{J}(.809, -.588) %cos324,sin324
\autonodetext{A}[e]{Mr. Landau} \autonodetext{B}[e]{8}
\autonodetext{C}[n]{7} \autonodetext{D}[n]{6}
\autonodetext{E}[n]{5} \autonodetext{F}[w]{4}
\autonodetext{G}[s]{3} \autonodetext{H}[s]{2}
\autonodetext{I}[s]{1} \autonodetext{J}[s]{0} \edge{B}{A}
\edge{B}{C} \edge{B}{D} \edge{B}{E} \edge{B}{F} \edge{B}{G}
\edge{B}{H} \edge{B}{I} \edge{C}{D} \edge{C}{E} \edge{C}{F}
\edge{C}{G}\edge{C}{H} \edge{D}{E} \edge{D}{F} \edge{D}{G}
\edge{E}{F}\edge{C}{A}\edge{D}{A} \edge{E}{A}
\end{graph}\vspace*{.5in}\caption{Example
\ref{exa:landaus}.} \label{fig:landaus3}
\end{minipage}
\end{figure}

% Not clear that this is important in our context.
% CAC, 6/19/14.
%\section{Graphic Sequences}
%
%\begin{df}
%A sequence of non-negative integers is {\em graphic} if there
%exists a graph whose degree sequence is precisely that sequence.
%\end{df}
%
%\begin{exa}
%The sequence $1,1,1$ is graphic, since $K_3$ is a graph with this
%degree sequence, and in general, so is the sequence
%$\underbrace{n,n,\ldots , n}_{n+1\ n\mathrm{'s}}$, since $K_{n+1}$
%has this degree sequence. The degree sequence
%$1,\underbrace{2,2,\ldots , 2}_{n\ \mathrm{twos}}, 1$ is graphic,
%since $P_{n+1}$ has this sequence. The degree sequence
%$\underbrace{2,2,\ldots , 2}_{n\ \mathrm{twos}}$ is graphic, since
%$C_{n}$ has this sequence. The
%sequence $0,1,2,3,4,5,6,7,8$ is graphic, but the sequence
%$1,2,3,4,5,6,7,8,9$ is not according to Example \ref{exa:landaus}.
%\end{exa}
%
%\begin{exer}
%Prove, using induction, that the sequence
%$$n, n, n - 1, n - 1, \ldots , 4, 4, 3, 3, 2, 2, 1, 1$$
%is always graphic. 
%\begin{answer} 
%The sequence $1, 1$ is clearly graphic.
%Assume that the sequence
%$$n - 1, n - 1, \ldots , 4, 4, 3, 3, 2, 2, 1, 1$$
%is graphic and
%add two vertices, $u, v.$ Join $v$ to one vertex of degree $n -
%1$, one of degree of $n - 2,$, etc., one vertex of degree 1. Since
%$v$ is joined to $n - 1$ vertices, and $u$ so far is not joined to
%any vertex, we have a sequence
%$$n, n - 1,  n - 1, n - 1,  n - 2, n - 2, \ldots , 4, 4, 3, 3, 2, 2, 1, 0.$$
%Finally, join $u$ to $v$ to obtain the sequence
%$$n, n, n - 1, n - 1, \ldots , 4, 4, 3, 3, 2, 2, 1, 1.$$
%\end{answer}
%\end{exer}
%
%
%\begin{exer} Seven friends go on holidays. They decide that each will
%send a postcard to three of the others. Is it possible that every
%student receives postcards from precisely the three to whom he
%sent postcards? Prove your answer!
%\begin{answer}
%The sequence $3, 3, 3, 3, 3, 3, 3$ is not graphic, as the number
%of vertices of odd degree is odd. Thus the given condition is not
%realizable.
%\end{answer}
%\end{exer}



% Not relevant in our context.
% CAC, 8/8/13
%
%\begin{figure}[h]
%\begin{minipage}{3cm}
%\rput(2,0){\begin{graph}(1,1)
% \roundnode{A}(-1,0)
%  \roundnode{B}(0,0)
%   \roundnode{C}(1,0)
%   \autonodetext{A}[w]{$A$}
%\autonodetext{B}[e]{$B_i$}
% \autonodetext{C}[e]{$D_j$}
% \bow{A}{B}{0.5}[\graphlinedash{3}]
%  \bow{A}{C}{-0.5}
%   \end{graph}}
%\vspace*{.5in}\caption{Theorem
%\ref{thm:Havel_Hakimi}.}\label{fig:hvhk1}
%\end{minipage}
%\hfill
%\begin{minipage}{3cm}
%\rput(2,0){\begin{graph}(1,1) \roundnode{A}(-1,0)
%  \roundnode{B}(0,0)
%   \roundnode{C}(1,0)
%   \autonodetext{A}[w]{$A$}
%\autonodetext{B}[e]{$C_j$}
% \autonodetext{C}[e]{$B_i$}
% \bow{A}{B}{0.5}
%  \bow{A}{C}{-0.5}[\graphlinedash{3}]
%\end{graph}}
%\vspace*{.5in}\caption{Theorem
%\ref{thm:Havel_Hakimi}.}\label{fig:hvhk2}
%\end{minipage}
%\hfill
%\begin{minipage}{3cm}
% \rput(2,0){\begin{graph}(1,1)
%\roundnode{A}(-1,0) \roundnode{B}(-.5,0)
% \roundnode{C}(.5,0)
%  \roundnode{D}(1,0)
%\autonodetext{A}[w]{$A$} \autonodetext{B}[s]{$B_i$}
%\autonodetext{C}[s]{$D$} \autonodetext{D}[e]{$C_j$}
%\bow{A}{B}{.5}[\graphlinedash{1}] \bow{B}{C}{.5}
%\bow{C}{D}{.5}[\graphlinedash{1}]\bow{D}{A}{.5}
% \end{graph}}
%\vspace*{.5in}\caption{Theorem
%\ref{thm:Havel_Hakimi}.}\label{fig:hvhk3}
%\end{minipage}
%\hfill
%\begin{minipage}{3cm}
% \rput(2,0){\begin{graph}(1,1)
% \roundnode{A}(-1,0) \roundnode{B}(-.5,0)
% \roundnode{C}(.5,0)
%  \roundnode{D}(1,0)
%\autonodetext{A}[w]{$A$} \autonodetext{B}[s]{$B_i$}
%\autonodetext{C}[s]{$D$} \autonodetext{D}[e]{$C_j$} \bow{A}{B}{.5}
%\bow{B}{C}{.5}[\graphlinedash{1}]
%\bow{C}{D}{.5}\bow{D}{A}{.5}[\graphlinedash{1}]
%\end{graph}}
%\vspace*{.5in}\caption{Theorem
%\ref{thm:Havel_Hakimi}.}\label{fig:hvhk4}
%\end{minipage}
%\end{figure}
%\begin{thm}[Havel-Hakimi]\label{thm:Havel_Hakimi}
%The two degree sequences $$I: \qquad  a \geq b_1 \geq b_2 \geq
%\cdots \geq b_a \geq c_1 \geq c_2 \geq \cdots \geq  c_n,$$
%$$II: \qquad b_1-1, b_2-1,  \cdots ,
%b_a-1, c_1,  c_2, \cdots , c_n,$$ are simultaneously graphic.
%\end{thm}
%\begin{pf}
%Assume first that the sequence $II$ is graphic. There is a graph
%$G'$ with degree sequence equal to sequence $II$. We construct the
%graph $G$ from $G'$ by adding a vertex and connecting it to the
%vertices whose degrees are $ b_1-1, b_2-1,  \cdots , b_a-1$. Then
%$G$ is a graph whose degree sequence is sequence $I$, and so
%$II\implies I$.
%
%\bigskip
%
%Assume now that sequence $I$ is graphic. Let $A, B_i, C_i$ be
%vertices with $\deg A = a, \deg B_i = b_i$, and $\deg C_i = c_i$,
%respectively. If $A$ were adjacent to all the $B_i$, our task is
%finished by simply removing $A$. So assume that there is $B_i$ to
%which $A$ is not adjacent, and a $C_j$ to which $A$ is adjacent.
%As the sequence is arranged in decreasing order, we must have $b_i
%\geq c_j$. If it happens that $b_i = c_j$, we then simply exchange
%$B_i$ and $D_j$ (see figures \ref{fig:hvhk1} and \ref{fig:hvhk2}).
%If $b_i>c_j$ then $B_i$ has at least one more neighbour than
%$C_j$. Call this neighbour $D$. In this case we remove the edges
%$AC_j$ and $B_iD$ and add the edges $AB_i$ and $DC_j$ to obtain a
%new graph with the same degree sequence as $II$. See figures
%\ref{fig:hvhk3} and \ref{fig:hvhk4}. This process is iterated
%until $A$ is adjacent to all the $B_i$. This finishes the proof.
%\end{pf}
%
%
%\begin{exa}
%Determine whether the degree sequence $6, 5, 4, 3, 2, 2, 2, 2$ is
%graphic. \end{exa}Solution: Using the Havel-Hakimi Theorem
%successively we have
%$$6, 5, 4, 3, 2, 2, 2,
%2 \rightarrow$$
%$$4, 3, 2, 1, 1, 1, 2 \rightarrow$$
%$$4, 3, 2, 2, 1, 1, 1\rightarrow$$
%$$2, 1, 1, 0, 1, 1\rightarrow$$
%$$2, 1, 1, 1, 1, 0\rightarrow$$
%$$0, 0, 1, 1, 0\rightarrow$$
%$$1, 1, 0, 0, 0.$$This last sequence is graphic. By the Havel-Hakimi Theorem, 
%the original sequence is
%graphic.

\newpage

%\section{Traversability}\label{sec:traverse}

\subsection{Trails, Paths, and Cycles}\label{sec:trails}
\begin{df}
\index{trail}
Recall that a {\bf trail} is a walk where all the edges are
distinct. 
An {\bf Eulerian trail}\index{trail!Eulerian}\index{Eulerian Trail}  
on a graph $G$is a trail that traverses every edge of
$G$. A {\bf tour}\index{tour} of $G$ is a closed walk that
traverses each edge of $G$ at least once. An {\bf Euler tour} 
(or {\bf Euler cycle})\index{tour!Euler}\index{Euler tour}\index{Euler cycle}  
on $G$ is a tour traversing each edge of $G$ exactly once, that is, a
closed Euler trail. 
A graph is {\bf Eulerian}\index{Eulerian graph}\index{graph!Eulerian} if it contains an Euler tour.
\end{df}

It turns out there is a very easy way to determine whether or not a graph has an Euler tour.

\begin{thm}\label{thm:Eulerian}
A nonempty connected graph is Eulerian if and only if it has no
vertices of odd degree.
\begin{pf}
Assume first that $G$ is Eulerian, and let $C$ be an Euler tour of
$G$ starting and ending at vertex $u$. Each time a vertex $v$   is
encountered along $C$, two of the edges incident to $v$ are
accounted for. Since  $C$ contains every edge of $G$, $d(v)$ is
then even for all $v\neq u$. Also, since $C$ begins and ends in
$u$, $d(u)$ must also be even.

Conversely, assume that $G$ is a connected nonEulerian graph with
at least one edge and no vertices of odd degree. Let $W$ be the
longest walk in $G$ that traverses every edge at most once: $$W =
v_0, v_0v_1, v_1, v_1v_2, v_2, . . . , v_{n-1}, v_{n-1}v_n,
v_n.$$ Then $W$ must traverse every edge incident to $v_n$,
otherwise, $W$ could be extended into a longer walk. In
particular, $W$ traverses two of these edges each time it passes
through $v_n$ and traverses $v_{n-1}v_n$ at the end of the walk.
This accounts for an odd number of edges, but the degree of $v_n$
is even by assumption. Hence, $W$ must also begin at $v_n$, that
is, $v_0 = v_n$. If $W$ were not an Euler tour,  we could find an
edge not in $W$ but incident to some vertex in $W$ since $G$ is
connected. Call this edge $uv_i$. But then we can construct a
longer walk:
$$u, uv_i, v_i, v_iv_{i+1}, . . . , v_{n-1}v_n, v_n, v_0v_1, .
. . , v_{i-1}v_i, v_i.$$ This contradicts the definition of $W$,
so $W$ must be an Euler tour.
\end{pf}
\end{thm}


The following problem is perhaps the originator of graph theory.

\begin{exa}[K\"{o}nigsberg Bridge Problem] \label{exa:koningsberg}\index{K\"{o}nigsberg Bridge}
The town of K\"{o}nigsberg (now called
Kaliningrad) was built on an island in the Pregel River. The
island  sat near where two branches of the river join, and the
borders of the town spread over to the banks of the river as
well as a nearby promontory. Between these four land masses, seven
bridges had been erected. The townsfolk used to amuse themselves
by crossing over the bridges and asked whether it was possible to
find a trail starting and ending in the same location allowing one
to traverse each of the bridges exactly once. Figure
\ref{fig:koningsberg} has a graph-theoretic model of the town,
with the seven edges of the graph representing the seven bridges.
By Theorem \ref{thm:Eulerian}, this graph is not Eulerian so it is
impossible to find a trail as the townsfolk asked.
\end{exa}

\begin{figure}[ht]
\centering
\begin{pspicture}(0,0)(2,2.5)
\rput(2,2){\begin{graph}(1,1) \roundnode{A}(0,1)
\roundnode{B}(0,0)\roundnode{C}(0,-1)
\roundnode{D}(2,0)\autonodetext{A}[n]{A}\autonodetext{B}[w]{B}\autonodetext{C}[s]{C}\autonodetext{D}[e]{D}
\edge{A}{D}\edge{B}{D}\edge{C}{D}\bow{A}{B}{.1}\bow{B}{C}{.1}\bow{C}{B}{.1}\bow{B}{A}{.1}
\end{graph}}
\end{pspicture}
\caption{Model of the bridges in K\"{o}nigsberg from Example~\ref{exa:koningsberg}.}\label{fig:koningsberg}
\end{figure}

\begin{df}
\index{Hamiltonian cycle}
\index{Hamiltonian graph} \index{graph!Hamiltonian}
A
{\bf Hamiltonian cycle} in a graph is  a cycle passing through
every vertex.  $G$ is {\bf Hamiltonian} if it contains a
Hamiltonian cycle.
\end{df}
Unlike Theorem \ref{thm:Eulerian}, there is no simple
characterization of all graphs with a Hamiltonian cycle. 
In fact, the problem of determining whether or not a graph contains a Hamiltonian cycle
is one of the most famous {\em NP-Complete} problems.
The details are beyond the scope of this book, but briefly (and oversimplifying a bit), 
{\em NP-Complete}  is a class of problems that are all equivalent in
the sense that if any of them can be solved in polynomial time,
then they can all be solved in polynomial time. Further, nobody currently knows whether or not
any of them can be solved in polynomial time. This leads to the so-called {\em P versus NP} problem,
one of the most important open problems in theoretical computer science. 
(Again, the details of precisely what this means are beyond the scope of this book.)

Coming back to the Hamiltonian cycle problem, we do have the following one-way result.

\begin{thm}[Dirac's Theorem, 1952]\label{thm:dirac}
\index{Dirac's Theorem} Let $G= (V,E)$ be a graph with $n =
|V| \geq 3$ vertices where each vertex has degree $\geq
\frac{n}{2}$. Then $G$ is Hamiltonian.
\begin{pf}
Arguing by contradiction, suppose $G$ is a maximal non-Hamiltonian
graph with $n\geq 3$, and that $G$ has more than $3$ vertices. Then
$G$ cannot be complete. Let $a$ and $b$ be two non-adjacent
vertices of $G$. By definition of $G$, $G+ab$ is Hamiltonian, and
each of its Hamiltonian cycles must contain the edge $ab$. Hence,
there is a Hamiltonian path $v_1v_2\ldots v_n$ in $G$ beginning at
$v_1=a$ and ending at $v_n = b$. Put $$S =\{v_i: av_{i+1}\in E\}
\qquad \mathrm{and}\qquad \{v_j:v_jb\in E\}.
$$As $v_n\in S\cap T$, we must have $|S\cup T| = n$.
Moreover, $S\cap T = \varnothing$, since if $v_i\S \cap T$ then
$G$ would have the Hamiltonian cycle
 $$v_1v_2\cdots v_iv_nv_{n-1}\cdots v_{i+1}v_1,$$
 as in the following figure, contrary to the assumption that $G$ is
 non-Hamiltonian.
 
\bigskip

\rput(6,0){\begin{graph}(1,1)
\bow{E}{A}{.2}[\graphlinewidth*{2.5}]\bow{D}{G}{.2}[\graphlinewidth*{2.5}]
\roundnode{A}(-5,0)\autonodetext{A}[s]{$v_1$}\roundnode{B}(-4,0)\autonodetext{B}[s]{$v_2$}
\roundnode{C}(-3,0)\autonodetext{C}[s]{$v_2$}\roundnode{D}(-1,0)\roundnode{E}(0,0)\roundnode{F}(4,0)
\roundnode{G}(5,0) \autonodetext{D}[s]{$v_i$}
\autonodetext{E}[s]{$v_{i+1}$} \autonodetext{F}[s]{$v_{n-1}$}
\autonodetext{G}[s]{$v_{n}$}
\edge{A}{B}[\graphlinewidth*{2.5}]\edge{B}{C}[\graphlinewidth*{2.5}]\edge{C}{D}[\graphlinedash{3}
\graphlinewidth*{2.5}] \edge{D}{E}\edge{E}{F}[\graphlinedash{3}
\graphlinewidth*{2.5}]\edge{F}{G}[\graphlinewidth*{2.5}]\end{graph}}
\vspace{.6in}

But then $$d(a) + d(b) = |S| + |T| = |S\cup T| + |S\cap T| < n.$$
But since we are assuming that $d(a)\geq \dfrac{n}{2}$ and $d(b)\geq
 \dfrac{n}{2}$, we have arrived at a contradiction.
\end{pf}
\end{thm}


%\newpage

\subsection{Planar Graphs}\label{sec:planar}
\begin{df}
A graph is {\bf planar}\index{planar graph}\index{graph!planar} 
if it can be drawn in
a plane with no intersecting edges.
Such a drawing is called a {\bf planar embedding} of the graph.
\end{df}
\begin{exa}
Although the usual way $K_4$ is drawn has two edges intersect, 
it is planar as shown in figure~\ref{fig:faces_k4}.
It is important to understand that being planar means you {\em can} draw it with no intersecting edges, not that every way of drawing it has no edges intersecting.
\end{exa}
\begin{figure}[h]
\centering
\begin{graph}(2,1)
\roundnode{A}(-1,1)\autonodetext{A}[w]{A}
\roundnode{B}(1,1)\autonodetext{B}[ne]{B}
\roundnode{C}(1,-1)\autonodetext{C}[e]{C}
\roundnode{D}(-1,-1)\autonodetext{D}[w]{D}
\edge{A}{B}\edge{A}{C}\edge{B}{C}\edge{C}{D} \edge{D}{A}
\bow{B}{D}{.8} \freetext(1.5,0){{\bf 2}}\freetext(.5,.5){{\bf 3}}
\freetext(-.5,-.5){{\bf 4}}\freetext(2.5,.0){{\bf 1}}
\end{graph}
\vspace*{.75in}\caption{A planar embedding of $K_4$.}\label{fig:faces_k4}
\end{figure}

\begin{exer}
Draw a planar embedding of $K_4$ that does not have curved edges.

\vspace*{1.25in}
\begin{answer}
Did you draw a triangle with a vertex in the middle connected to the
three vertices of the triangle? I thought so!
\end{answer}
\end{exer}

\begin{df}\index{face, planar graph}
A {\bf face}\index{face} of a planar graph is a region bounded by
the edges of the graph.
\end{df}

\begin{exa}\label{exa:faces_k4}
$K_4$ has $4$ faces, labeled 1 through 
4 in Figure~\ref{fig:faces_k4}. 
Face {\bf 1}, which extends indefinitely, 
is called the {\em outside face}.\index{outside face}
\end{exa}


Here are a few results about planar graphs.
These theorems use $v$ and $e$ instead of $n$ and $m$ because although computer scientists
often use $n$ and $m$, graph theorists seem to prefer $v$ and $e$. And you should get used to
the fact that not everybody uses the same notation, so it's good for you to see different letters used.

\begin{thm}[Euler's Formula] \index{Euler's formula}
For every drawing of a connected
planar graph with $v$ vertices, $e$ edges, and $f$ faces the
following formula holds: $$v-e+f = 2.  $$
\begin{pf}
The proof is by induction on $e$. Let $P(e)$ be the proposition
that $v - e + f = 2$ for every drawing of a graph $G$ with $e$
edges. If  $e = 0$ and it is connected, then we must have  $v = 1$
 and hence $f = 1$, since there is only the outside face. Therefore,
 $v - e + f = 1 - 0 + 1 = 2$, establishing $P(0)$
 
Assume now $P(e)$ is true, and consider a connected graph $G$
with $e + 1$ edges. Either
\begin{dingautolist}{202}
\item  $G$ has no cycles. Then there is only the outside face, and
 so $f=1$. Since there are $e+1$ edges and $G$ is connected, we must have $v=e+2$. This gives
 $(e+2)-(e+1)+1 = 2-1+1 = 2$, establishing $P(e+1)$.
\item or $G$ has at least one cycle. Consider a spanning
tree of $G$ and an edge $uv$ in the cycle, but not in the tree.
Such an edge is guaranteed by the fact that a tree has no cycles.
Deleting $uv$ merges the two faces on either side of the edge and
leaves a graph $G'$ with only $e$ edges, $v$ vertices, and $f$
faces. $G'$ is connected since there is a path between every pair
of vertices within the spanning tree. So $v - e + f = 2$ by the
induction assumption $P(e)$. But then
$$v - e + f
= 2 \implies (v) - (e+1) + (f+1) = 2 \implies v-e+f = 2,
$$establishing $P(e+1)$.
\end{dingautolist}
This finishes the proof.
\end{pf}\end{thm}

\begin{thm}
\begin{lenum}
\item 
Every simple planar graph with $v\geq 3$ vertices has  $e \leq
3v-6$ edges. 
\item 
Every simple planar graph with $v\geq 3$ vertices and
which does not have $C_3$ as a subgraph has $e\leq 2v-4$ edges.
\end{lenum}
\label{thm:kurat} % This is here because when I put it in other places it linked to page 0 for some unknown reason.
\begin{pf}
If $v = 3$, both statements are plainly true so assume that  $G$
is a maximal planar graph with $v\geq 4$. We may also assume that
$G$ is connected, otherwise, we may add an edge to $G$. Since $G$
is simple, every face has at least $3$ edges in its boundary. If
there are $f$ faces, let $F_k$ denote the number of edges on the
$k$-th face, for $1 \leq k \leq f$. We then have $$F_1 + F_2
\cdots + F_f \geq 3f.$$ Also, every edge lies in the boundary of
at most two faces. Hence if $E_j$ denotes the number of faces that
the $j$-th edge has, then $$2e  \geq E_1 + E_2 + \cdots + E_e.
$$Since $E_1 + E_2 + \cdots + E_e = F_1 + F_2
\cdots + F_f$, we deduce that $2e \geq 3f$. By Euler's Formula we
then have $e\leq 3v -6$.

\medskip
The second statement follows for $v=4$ by inspecting all graphs
$G$ with $v = 4$. Assume then that $v\geq 5$ and that $G$ has no
cycle of length $3$. Then each face has at least four edges on its
boundary. This gives $2e \geq 4f$ and by Euler's Formula, $e \leq
2v -4$.
\end{pf}
\end{thm}

To be clear, Theorem~\ref{thm:kurat} part (a) implies that a
graph with at least 3 vertices and more than $3v-6$ edges cannot be planar (the contrapositive of the statement).
Similarly for part (b).

\begin{exa}
$K_5$ is not planar by Theorem~\ref{thm:kurat} since
$K_5$ has $\binom{5}{2} = 10$ edges and $10 > 9 = 3(5)-6$.
\end{exa}

\begin{exer}
Prove that $K_{3,3}$ is not planar.

\noindent
\answerlinetwo

\begin{answer}
Notice that $K_{3,3}$ does not have $C_3$ as a subgraph. 
Since $K_{3,3}$ has $3\cdot 3 = 9$ edges and $9 > 8 = 2(6)-4$,
Theorem~\ref{thm:kurat} part (b) implies that $K_{3,3}$ is not planar.
\end{answer}
\end{exer}

% Interesting, but not terribly relevant.
%  CAC, 6/19/14.
%\begin{df}
%A {\bf polyhedron} is a convex, three-dimensional region bounded
%by a finite number of polygonal faces.
%\end{df}
%\begin{df}
%A {\bf Platonic solid} is a polyhedron having congruent regular
%polygon as faces and having the same number of edges meeting at
%each corner.
%\end{df}
%By puncturing a face of a polyhedron and spreading its surface
%into the plane, we obtain a planar graph.
%\begin{exa}[Platonic Solid Problem]\index{solid!Platonic} How many
%Platonic solids are there? If $m$ is the number of faces that meet
%at each corner of a polyhedron, and $n$ is the number of sides on
%each face, then, in the corresponding planar graph, there are $m$
%edges incident to each of the $v$ vertices. As each edge is
%incident to two vertices, we have $mv = 2e$, and if each face is
%bounded by $n$ edges, we also have $nf = 2e$. It follows from
%Euler's Formula that
%$$\dfrac{2e}{m} - e + \dfrac{2e}{n} = 2 \implies \dfrac{1}{m} +
%\dfrac{1}{n} = \dfrac{1}{e} + \dfrac{1}{2}.$$ We must have $n \geq
%3$ and $m \geq 3$ for a nondegenerate polygon. Moreover, if either
%$n$ or $m$ were $\geq 6$  then  $$\leq \dfrac{1}{3} + \dfrac{1}{6}
%= \dfrac{1}{2} < \dfrac{1}{e} + \dfrac{1}{2}.$$Thus we only need
%to check the finitely many cases with $3 \leq n, m \leq 5$. The
%table below gives the existing polyhedra.
%$$\begin{array}{|ll|lll|l|}\hline n & m & v &  e &
%f & \mathrm{polyhedron} \\
%\hline 3&  3 &  4 &  6 &  4 & \mathrm{tetrahedron} \\
%\hline 4 &  3 &  8 &  12 &  6 & \mathrm{cube} \\
%\hline 3 &  4 &  6 &  12 &  8 & \mathrm{octahedron} \\
%\hline 3 & 5 & 12 & 30&  20 & \mathrm{icosahedron} \\
%\hline 5 & 3 & 20 & 30 & 12 & \mathrm{dodecahedron} \\
%\hline
%\end{array}$$
%\end{exa}
%
%\begin{exa}[Regions in a Circle]Prove that the chords determined
%by $n$ points on a circle cut the interior into $1 + \binom{n}{2}
%+ \binom{n}{4}$ regions provided no three chords have a common
%intersection.
%\begin{sol}By viewing  the points on the circle and the
%intersection of two chords as vertices, we obtain a plane graph.
%Each intersection of the chords is determined by four points on
%the circle, and hence our graph has $v = \binom{n}{4} + n$
%vertices. Since each vertex inside the circle has degree $4$ and
%each vertex on the circumference of the circle has degree $n+1$,
%the Handshake Lemma (Theorem \ref{thm:handshake_lemma}) we have a
%total of $$e = \dfrac{1}{2}\left(4\binom{n}{4} + n(n+1)\right)
%$$edges. Discounting the outside face, 
%our graph has 
%$$f-1 = 1 + e-v = 1 + 2\binom{n}{4} + \dfrac{n^2}{2} + \dfrac{n}{2}
%- \left( \binom{n}{4} + n\right) = 1 + \binom{n}{2} + \binom{n}{4}
%$$faces or regions.
%\end{sol}
%\end{exa} 

%\newpage
%
%\section{Problems}
%
%{\Huge Add some problems here!}




\newpage

\subsection{Minimum Spanning Trees}\label{sec:mst}

Recall that a {\bf spanning tree} of $G$ is a subgraph
$T$ of $G$ which is a tree that spans $G$. 
In other words, it contains all of the vertices of $G$.
Some problems can be boiled down to trying to find a
spanning tree of a weighted graph such that the sum
of the weights of all of the edges of the spanning tree 
is as small as possible. This section formally
defines this problem and then presents two algorithms 
to solve it.

\begin{df}\index{weight, tree}
If $T$ is a tree, the {\bf weight} of a $T$ is the sum of the
weights of its edges.  That is, 
$$w(T)=\sum_{(u,v)\in T} w(u,v).$$
\end{df}

\begin{df}\index{minimum spanning tree}\index{MST}
Let $G=(V,E)$ be a connected, weighted graph.
A {\bf minimum spanning tree (MST)} of $G$ is spanning tree 
$T$ of minimum weight.
\end{df}

It should be clear that a minimum spanning tree always exists.

\begin{exa}
\begin{multicols}{2}
Here is an example of a graph $G$ along with 
two different minimum spanning trees of $G$,
$T_1$ and $T_2$.
Notice that $w(T_1)=w(T_2)=18$, and try as you might,
you will be unable to find a spanning tree of weight less than 18.
\\

\ \ \ \sputfig{graphs/fig/mst1.eps}{.5}
\end{multicols}
\end{exa}


Minimum spanning trees can be constructed in a greedy fashion.
There are two common algorithms to construct MSTs, 
{\bf Kruskal's algorithm} and {\bf Prim's algorithm}.
\index{Kruskal's algorithm}\index{Prim's algorithm}
Both of these algorithms use the same basic ideas, but
in a slightly different fashion.
In the rest of this section
we will consider a general approach to find MSTs,
prove that the general approach works,
and show how to implement 
Kruskal's  and Prim's algorithms based on the approach.

\begin{df}\index{cut}\index{cross edge}\index{light edge}
Let $G=(V,E)$ be a connected, weighted graph.
\begin{itemize}
\item
A {\bf cut} $(S,V-S)$ of $G$ is a partition of the vertices $V$.
\item
An edge $(u,v)\in E$ is said to {\bf cross} the cut if 
one of the endpoints is in $S$, and the other is in $V-S$.
\item
The set of edges which cross a cut are the {\bf cross edges}.
\item
A cut {\bf respects} a set $A$ of edges if $A$
does not contain any cross edges.
\item
A cross edge of minimum weight is 
called a {\bf light edge}.
\end{itemize}
\end{df}

\begin{exa}\label{exa:cuts}
Here is an example of some of the terminology applied to graph.\\

\sputfig{graphs/fig/mst2.eps}{.6}
\end{exa}


The generic MST algorithm can be described as follows.
Let $A$ be the edges a minimal spanning tree of $G$.
The MST algorithm ``grows'' the spanning tree 
one edge at a time.
It starts with set $A=\emptyset$, which is clearly a subset of
every minimum spanning tree.
At each step, the algorithm adds an edge 
$(u,v)$ to $A$ so that the set $A\cup\{(u,v)\}$ is a
subset of some minimum spanning tree.
Such an edge $(u,v)$ is called a {\bf safe edge},
\index{safe edge} because
we can safely add it to the set $A$ and still continue.

The algorithm is simple:
\begin{proc}
The generic MST algorithm.
\begin{code}
MST(G)
 A=EmptyList
 While NOT IsSpanningTree(G,A) 
     e = SafeEdge(G,A)
     Insert(A,e)
 return A
\end{code}
\end{proc}

In some sense, this is not technically an algorithm. Why?
In order for it to be an algorithm, we need each step to be
clearly executable. But there are at least two steps that are
unclear. First, how does \verb+IsSpanningTree(G,A)+ determine whether
or not $A$ is a spanning tree of $G$?
Second, how does \verb+SafeEdge(G,A)+ find a safe edge?
Until we answer these questions we really just have a vague outline
of an algorithm. 

We will start by attempting to address the first question.
Consider $A$, the partial set of edges
that we are building into a spanning tree during the MST algorithm
Notice that the graph $G_A=(V,A)$ is a forest, with some of the
trees consisting of just a single node, 
and that each tree in the forest $G_A$ is a {\bf connected component}.
At every step of the algorithm, MST adds an edge to the set $A$.
which results in the merger of two trees into one.

\begin{exa}
Consider the graph below where the edges of $A$ are darkened
and we have labeled the trees of the forest $T_1$ through $T_5$.
If we add the dashed edge to $A$ in the next step, the result will
be that trees $T_3$ and $T_4$ will be merged into a single tree.\\

\ \ \ \ \ \ \sputfig{graphs/fig/mst4.eps}{.5}
\end{exa}

So what does this have to do with determining whether or not 
$A$ is a spanning tree? Simple: At the beginning, 
$|A|=0$, and we have $|V|$ trees in the forest, one for each vertex.
At each step of the algorithm, we merge two trees. Therefore, 
all we have to do is iterate through $|V|-1$ times and we will be left
with a single tree--a spanning tree!

Now the harder question: how do we find safe edges. But first,
a question that may have occurred to you: 
How do we know that there are any safe edges? 
The algorithm assumes they exist,
so we need to be certain that they always do. 
Luckily, we defined safe edges in such a way that that {\em have to}
exist as long as we follow the algorithm correctly. 
Let us argue this so it is more clear. 

At the beginning of the algorithm, $A=\emptyset$, and
any edge in any minimum spanning tree is safe. Since we know
at least one spanning tree exists, then any edge in that tree is safe. 
So we know we that at the beginning there are safe edges.

During each iteration, we add a safe edge to $A$. Recall that
by definition a safe edge is an edge 
$(u,v)$ such that the set $A\cup\{(u,v)\}$ is a
subset of some minimum spanning tree.
Thus, every time we add a safe edge, the revised set $A$ is
guaranteed to still be a subset of some minimum spanning tree $T$.
Thus, any edge from $T-A$ is a safe edge, and since the algorithm 
has not finished yet, there must be at least one safe edge (i.e.
some edge in $T$ that is not in $A$. If that sounds a bit convoluted,
it is not surprising. Read it slowly and carefully a few more times 
and I think you will eventually catch on.

Now we know there are safe edge, how do we find them?

Actually, it's not that hard to find safe edges, as
we will see next.

We begin with an important result.

\begin{thm}\label{thm:mst0}
Let $G\!=\!(V,E)$ be a connected, weighted graph, and 
	\begin{itemize}
	\item
	$A\subseteq E$ a subset of some MST for $G$,
	\item
	$(S,V-S)$ be any cut of $G$ that respects $A$, and
	\item
	$(u,v)$ be a light edge of $(S,V-S)$.
	\end{itemize}
Then the edge $(u,v)$ is safe for $A$.
\begin{pf}
The idea of the proof is as follows: 
Let $T$ be a MST of $G$ containing $A$, 
and $(u,v)$ an edge as described above.  
If $(u,v)\in T$, then $(u,v)$ is safe for $A$, and we are done.
If $(u,v)\not\in T$, we need to find an MST $T^{\prime}$ such that 
$A\cup\{(u,v)\}\subseteq T^{\prime}$ (so that $(u,v)$ is safe).
In order to do this, we will find an edge $(x,y)\in T$ such that the tree
$T^{\prime}=(T-\{(x,y)\})\cup\{(u,v)\}$ is an MST for $G$ that
contains $A$ and $(u,v)$.
that will mean that $(u,v)$ is safe for $A$.

The only step of the proof that is not clear is how to find
an edge $(x,y)\in T$ that will do as we desire. 
To see how we can do that,
consider the following diagram of the situation.

\ \ \ \ \ \ \sputfig{graphs/fig/mst3.eps}{.5}

Notice that the graph $T\cup\{(u,v)\}$ contains a cycle
(since adding any edge to a tree will create a cycle).
Since $(u,v)$ is a cross edge on the cycle, there must
be another cross edge on the cycle.
Let $(x,y)$ be such an edge.
We claim that 
$T^{\prime}=(T-\{(x,y)\})\cup\{(u,v)\}$ is an MST for $G$
containing $A$, so that $(u,v)$ is safe for $A$.
The edge $(x,y)$ is not in $A$, because the cut respects $A$.
Thus, $A$ is a subset of $T^{\prime}$.
Now all we need to show is that $T^{\prime}$ is an MST for $G$.

Proof that $T^{\prime}$ is an MST of $G$: 
Since $(x,y)$ is a cross edge and $(u,v)$ is a 
light edge,
$w(u,v)\leq w(x,y)$, which implies $w(u,v) - w(x,y)\leq 0$.
Therefore,  
$$w(T^{\prime})=w(T) +  w(u,v) - w(x,y)\leq w(T).$$
Since $T$ is an MST, $w(T)\leq w(T^{\prime})$.
Thus, $w(T)=w(T^{\prime})$, and $T^{\prime}$ is an MST for $G$.

To summarize, we have found a tree $T^{\prime}$ such that
$T^{\prime}$ is an MST of $G$,
$A$ is a subset of $T^{\prime}$,
$(u,v)\in T^{\prime}$, and $(u,v)\not\in A$, so 
$(u,v)$ is safe for $A$.
\end{pf}
\end{thm}


We can use Theorem~\ref{thm:mst0}
to prove the following.

\begin{thm}\label{thm:mst1}
$G\!=\!(V,E)$ be a connected, weighted graph, and
\begin{itemize}
\item
$A\subseteq E$ a subset of some MST for $G$, 
\item
$C$ be the edges in a connected component of $G_A\!=\!(V,A)$, and
\item
$(u,v)$ be a light edge of the cut $(C,V-C)$.
\end{itemize}
Then $(u,v)$ is safe for $A$.
\begin{pf}
Since the cut $(C,V-C)$ respects $A$, this follows from Theorem~\ref{thm:mst0}.
\end{pf}
\end{thm}

\begin{ques}
In the previous theorem, why does the cut $(C,V-C)$ respects $A$?
\vspace*{.75in}
\begin{answer}
Since $C$ is the set of edges of a connected component of
$G_A\!=\!(V,A)$,
then $A$ has no edges between vertices 
of $C$ and vertices of $V-C$.

Another way to think about it: If the cut $(C,V-C)$ does not respect $A$,
then there is some $(u,v)\in A$ such that $u\in C$ and $v\in V-C$.
But $C$ is the set of edges of a {\bf\em connected} component, 
and since $(u,v)\in A$, $v\in C$, contradicting the fact that $v\in V-C$.
Therefore, the cut $(C,V-C)$ respects $A$.
\end{answer}
\end{ques}

Theorem~\ref{thm:mst1} basically says that if $C$ is a subtree
of an MST, and $(u,v)$ is an edge of minimum weight 
with exactly one endpoint incident with $C$,
then $C\cup\{(u,v)\}$ is a subtree
of an MST for $G$.
Here are two ideas based on this theorem:
\begin{itemize}
\item {\bf Idea 1:}
Let $u$ be a vertex of $G$, and $(u,v)$ an edge
of minimum weight incident with $u$.
Then $(u,v)$ is contained in some MST of $G$.
(Here, $C=\{u\}$.)
\item {\bf Idea 2:}
If $(u,v)$ is an edge of minimal weight in $G$,
then $(u,v)$ is contained in some MST of $G$.
(Again, $C=\{u\}$.)
\end{itemize}

Next we provide high level descriptions of Prim's and Kruskal's algorithm
so you can hopefully see how they are each applying the ideas above.
Then we will give more detailed descriptions of both.

Prim's algorithm uses Theorem~\ref{thm:mst1} to build a single
tree into an MST by starting at some vertex and ``growing out''
from it.

\begin{exa}\index{Prim's algorithm}
Here is the idea of Prim's algorithm:
\begin{itemize}
\item
Pick some vertex $x$.
\item
Let $A=\{(x,y)\}$, where edge $(x,y)$ has minimum weight of 
edges incident with $x$.
\item
While $A$ is not an MST
\begin{itemize}
\item
Add to $A$ a minimum weight edge which has 
exactly one endpoint incident with $A$
\end{itemize}
\end{itemize}
\end{exa}

On the other hand, Kruskal's algorithm 
sorts the edges of the graph according to weight, and adds
edges (starting with the lightest) as long as they do not create
a cycle until a spanning tree is created.

\begin{exa}\index{Kruskal's algorithm}
The idea of Kruskal's Algorithm:
\begin{itemize}
\item
Let $A=\emptyset$.
\item
While A is not an MST
\begin{itemize}
\item
Add to $A$ a minimum weight edge that does not form a cycle.
\end{itemize}
\end{itemize}
\end{exa}

\newpage

\begin{ques}
Is Prim's Algorithm based on Idea 1 or 2? What about 
Kruskal's algorithm?
\vspace*{1in}
\begin{answer}
Since Idea 1 is considering adding edges adjacent to a given vertex,
Prim's algorithm is essentially using this idea, at least in the first step.
Although to be fair, Prim's algorithm is also essentially directly using
Theorem~\ref{thm:mst1}, where $C$ is the single tree it is growing.

On the other hand, Idea 2 is about using minimum weight edges 
without a specific endpoint in mind, which is what Kruskal's 
algorithm is doing.
\end{answer}
\end{ques}
Next, we dig into details for both of these algorithms. Let us
continue with Kruskal's algorithm.
Here is a slightly more detailed description of Kruskal's algorithm
that is adding a few important details that were missing. 
\begin{exa}
Kruskal's algorithm is as follows.
Given a graph $G=(V,E)$, treat $E$ as an array of edges.
\begin{itemize}
\item
Sort $E$ in ascending order.
\item
Set $A=\emptyset$
\item
For I = 1 to $|E|$\\
\hspace*{.1in} If $A\cup\{E[I]\}$ does not contain a cycle\\
\hspace*{.2in} $A=A\cup\{E[I]\}$
\item
Return $A$.
\end{itemize}
\end{exa}

We still need to answer an important questions in order to 
have a complete algorithm: 
When does $A\cup\{E[I]\}$ contains a cycle?
Notice that 
as the algorithm progresses, $A$ is a forest. 
Edges connecting two vertices in the same tree
will create a cycle.
Edges that go from one tree to another will not create a cycle.
So we will store each tree in a separate set.
Adding an edge connect two trees, so we merge the sets.
Based on these observations, we 
can now give a more detailed description of Kruskal's algorithm.

\begin{proc}\index{Kruskal's algorithm}
Kruskal's Algorithm

\begin{code}
Kruskal_MST(G)
   A=EmptySet
   ForAll v in V[G]
       Create_Set(v)
   SortAscending(E[G])
   ForAll edges e=(u,v) in E[G] //in sorted order
       If Set(u) != Set(v) 
          Insert(A,e)
          Set_Union(u,v)
   Return A
\end{code}
\end{proc}

There are still important unanswered questions that we will not answer, 
but we will ask: How do we implement the set data structure we need? 
It needs to store sets in a way that allows it to determine if
two elements are in the same set and create the union of two sets.
One solution is to use something called {\bf union-find algorithm}
(or {\bf union-find data structure}).
Unfortunately, we do not have time to go into the details of this
data structure.

Now that we have the algorithm, how good is it?
Let $n=|V|$ and $m=|E|$.
Creating the sets takes $O(n)$ time (constant time for each).
Sorting the edges takes $O(m\log m)$ time.
With a union-find data structure, 
it is possible to implement the sets so that 
comparing and unioning sets each take $O(\log n)$ time,
and adding an edge to $A$ takes constant time.
Since these operations are done at most $m$ times
(one for each edge), the cost of the \verb+ForAll+ loop
is $O(m\log n)$.
Therefore, the complexity of Kruskal's Algorithm is 
$O(m\log m)+O(m\log n)=O(m\log m)=O(m\log n)$.

\begin{ques}
Why is $O(m\log m)=O(m\log n)$?
\vspace*{1in}
\begin{answer}
Recall that the maximum number of edges a graph can have is 
$m\leq \binom{n}{2}=\frac{n(n-1)}{2}=O(n^2)$.
Therefore, $O(\log m)=O( \log n^2)=O(2\log n)=O(\log n)$.
Given this, it is clear that $O(m\log m)=O(m\log n)$
\end{answer}
\end{ques}


\begin{exa}
Here is an example of Kruskal's algorithm
(Follow down the columns here and then on the next page).
Note that we are not showing the details about the sets in this 
example, but following along like you would if you did it by hand.\\

\noindent
\sputfig{graphs/fig/kruskal1.eps}{.6} \hfill
\sputfig{graphs/fig/kruskal2.eps}{.6} 

\noindent
\sputfig{graphs/fig/kruskal3.eps}{.6} \hfill 
\sputfig{graphs/fig/kruskal4.eps}{.6}
\end{exa}

Before discussing Prim's algorithm, we first need a new data structure
called a {\em priority queue}.

\begin{df}\index{priority queue}
A {\bf priority queue} $Q$ is a data structure that supports the
following operations.
\begin{itemize}
\item
\verb+Q.insert(Item v,Key k)+ inserts the Item $v$ into $Q$ using the value of the Key $k$ to determine where it goes.
\item
\verb+Q.extractMin()+ removes and returns an Item from $Q$ 
that has a minimum key value. If there are multiple items with
the same minimum key value, which one gets returned depends
on the implementation. Most algorithms that
use priority queues do not care about which order tied keys are returned in.
\item
\verb+Q.decreaseKey(v,k)+ decreases the key value for Item $v$
to $k$ and moves it to the appropriate place in $Q$ 
based on the new key value.
\item
\verb+Q.insertAll(Item[] values,Key[] keys)+ 
inserts the Items from the list $values$ based on the corresponding 
$keys$ into $Q$.
\end{itemize}
\end{df}

There are multiple ways to implement a priority queue that has these
operations. A common choice, the {\em min heap}, allows us to
implement the priority queue so that the first three operations take 
$O(\log n)$ time if $Q$ contains $n$ items, and the last one
takes $O(n)$ time to insert a list of size $n$ into an empty $Q$.

Unlike Kruskal's algorithm, with Prim's algorithm
we grow a single tree $A$ into a minimum spanning tree.
An arbitrary vertex $r$ is picked, and the tree is grown
from that vertex.
At each step a light edge of the cut $(A,v-A)$ is added
to $A$. That is, we add an edge that connects a new vertex to $A$.
Since it adds a node not already connected to $A$, it remains a tree
since it is impossible to add a cycle if the new vertex is not 
already connected to some vertex in $A$.

In order to implement the algorithm, 
we need to have an efficient (greedy) way to determine
the light edge at each step.
To help facilitate this, for each node $x$, we will store
\begin{itemize}
\item
The $key(x)$ is the weight of the 
minimum weight edge that connects
$x$ to some vertex in $A$.
\item
The {\em predecessor} $p(x)$ is the vertex $y$ in $A$ such that $w(x,y)=key(x)$.
That is, of all of the nodes in $A$, $y$ is one such that
$w(x,y)$ is as small as possible.
\end{itemize}

We choose some root vertex $r$ to start growing our tree from 
and set $key(r)=0$, and $p(r)=NULL$.
Each node $x\neq r$ starts with $key(x)=\infty$
(or \verb+MAX_INT+ or whatever the largest value we can
store is).
As the algorithm progresses, the key and predecessor 
are continually updated to be the best possible.
We do not need to initialize $p(x)$ since it will eventually 
get set for every node (except $r$) by the algorithm.

When the algorithm is finished, the minimum spanning tree can
be constructed by using the values of $p$. For instance, 
if $p(a)=b$, then $(b,a)$ is an edge of the tree.

The value $key(x)$ only changes if some neighbor of $x$
is added to $A$.
When we add a node $y$ to $A$, we need to update the
$key$ values of the nodes adjacent to $y$.
Specifically, for each vertex $x$ adjacent to $y$, 
if $x$ is not already in $A$ and $w(x,y)<key(x)$, 
then we have found a cheaper way to 
connect $x$ to $A$, so we update the key and predecessor of $x$.

We will store the vertices in $V-A$ in a {\em priority queue} $Q$ 
based on $key(x)$.  
This allows us to pick the minimum weight edge
to add to $A$ at each step.

We are now ready for the full version of the algorithm.

\begin{proc}\index{Prim's algorithm}
Prim's Algorithm
\begin{code}
Prim_MST(G,r)
  PriorityQueue Q   // An empty priority queue
  Vertex[] p        // An array to store the predecessors
  int[] key         // An array to store the keys 
  ForAll u in G(V)  // Start all keys out at the largest possible value
      key[u]=Max_Int
  key[r]=0    // r has the lowest key so it get extracted first
  p[r]=NULL   // r will never have a predecessor since it is the root
  Q.insertAll(V,key) // Add all vertices to Q based on key values
  While NotEmpty(Q)     // Extract vertices one at a time until  
      u=Q.extractMin()  //    until we have a spanning tree   
      ForAll v adjacent to u  // See if we need to update neighbors
        if(v in Q and w(u,v) < key[v]) // If we found a cheaper edge:
            key[v]=w(u,v)              // update the key value
            Q.decreaseKey(v,key[v])    // update the location of v in Q
            p[v]=u                     // update the predecessor
\end{code}
\end{proc}


\begin{exa}
Here is an example of Prim's algorithm starting from root $r=a$.
Note: the diagram uses {\em nil} instead of {\em null} 
and $k$ instead of $key$ to save space.
The values of $k$ ($key$), $p$ and, $Q$ are shown at every step.
The blackened nodes are in $A$ and no longer in $Q$. Notice that
the values of $k$ for these nodes are also crossed out--that's just
to make it more apparent visually which nodes are no longer in $Q$.
\\

\sputfig{graphs/fig/prim1.eps}{.6}

The algorithm continues on the next page, but let's make 
a few observations before you continue.
The gray lines show the MST edges, which can be found by looking at $p(x)$ for vertices $x$ that are in $A$. For instance, on the third row
notice that $k(d)$ is crossed out and $p(d)=a$. That means
$(a,d)$ is an edge in the tree.
On the same line, $p(e)=d$, but that does not
mean $(d,e)$ is a tree edge because the algorithm has 
yet to confirm whether or not that is the cheapest way to connect
to $e$. In fact, as you will see on the next page, it is not.

\sputfig{graphs/fig/prim2.eps}{.6}
\end{exa}

\begin{exer}
Determine the computational complexity of \verb+Prim_MST+
given a graph with $n$ vertices and $m$ edges.
\vspace*{1.6in}
\begin{answer}
Lines 1-3 and  6-7 take constant time.
Lines 4-5 line takes $O(n)$ time since we are doing $n$ things
that take constant time.
Line 8 takes $O(n)$ time.
The \verb+While+ loop executes $n$ times,
each time calling \verb+Q.extractMin+ that takes $O(\log n)$
and the \verb+ForAll+ loop.
The code in the \verb+ForAll+ loop takes $O(\log n)$
time, and it executes at most $n-1$ times (since there can be
at most $n-1$ vertices adjacent to a given vertex).
Combining this together, the complexity is
$O(n)+O(n(\log n + (n-1)\log n))=O(n^2\log n)$ 
(Make sure you understand how I simplified this!).

Unfortunately, this actually overestimates the time complexity
of the algorithm. 
Obtaining a better bound takes a little thought and can be tricky
to understand, so read carefully!
Let's analyze the \verb+ForAll+ loop
more carefully. In particular, instead of counting up the operations
according to how the code is written (iterating over adjacency lists),
we will think about how many times the code inside that 
loop executes in total regardless of when it happens. 
Notice that the code in the \verb+ForAll+ loop
can execute at most two times for every edge--once from each
endpoint. If you cannot see why this is, definitely ask about it in class!
Thus, the complexity of the \verb+ForAll+ loop
is actually $O(m\log n)$ over all of the iterations of the
\verb+While+ loop. So the complexity of the \verb+While+
loop is $O(m\log n)$ plus the complexity of the \verb+Q.extractMin+ method called $n$ times, which we saw was $O(n\log n)$.
Given that, the complexity of the algorithm is 
$O(n)+O(m\log n)+O(n\log n)=O(m\log n)$.
\end{answer}
\end{exer}

Time for a bonus algorithm to solve a new problem called
the {\em single-source shortest path} problem.
\index{single-source shortest path}
Here we are given a weighted graph $G=(V,E)$ and a 
chosen vertex $r$, and we wish to know the shortest path
from $r$ to every other node in the graph. We will not
go into detail about this problem, but will simply present
an algorithm to solve it that will look very familiar.

\begin{proc}\index{Dijkstra's algorithm}
Dijkstra's Algorithm

Dijkstra's Algorithm is almost
identical to Prim's algorithm--the only differences
are marked below with \verb+<--+.
\begin{code}
Dijkstra(G,r)
  PriorityQueue Q
  Vertex[] p 
  int[] key
  ForAll u in G(V)
      key[u]=Max_Int
  key[r]=0
  p[r]=NULL
  Q.insertAll(V,key)
  While NotEmpty(Q)
      u=Q.extractMin()
      ForAll v adjacent to u
        if(v in Q and key[u]+w(u,v) < key[v]) <-- changed
            key[v]=key[u]+w(u,v)              <-- changed
            Q.decreaseKey(v,key[v]) 
            p[v]=u
\end{code}
\end{proc}

Part of learning about algorithms is learning how we can 
apply what we have learned to solve various problems 
to solve other related, and sometimes not-so-related, problems.
In this case, the minimum spanning tree and single source shortest path
are definitely not the same problem, but they both are attempting
to find spanning trees in the graph with certain properties.
When you realize that, it makes sense that we might be able to
slightly modify an algorithm for one problem to solve the other.



